\PassOptionsToPackage{dvipsnames,svgnames,table}{xcolor}
\documentclass[11pt]{article}
\usepackage[a4paper,margin=22mm]{geometry}
\usepackage[no-math]{fontspec}
\setmainfont{DejaVu Serif}
\usepackage{amsmath,amssymb,amsthm,mathtools,graphicx,xcolor,hyperref,xurl,xspace,ifthen,enumerate}
\usepackage[shortlabels]{enumitem}
\setlength{\emergencystretch}{4em}
\AtBeginDocument{\special{pdf:mapline pzdr Dingbats <uzdr.pfb}}
\SetMathAlphabet{\mathbf}{normal}{OT1}{cmr}{bx}{n}
\SetMathAlphabet{\mathbf}{bold}{OT1}{cmr}{bx}{n}
\newtheorem{theo}{Theorem}[section]
\newtheorem{lemm}[theo]{Lemma}
\newtheorem{thm}[theo]{Theorem}
\newtheorem{proposition}[theo]{Proposition}
\newtheorem{corollary}[theo]{Corollary}
\newtheorem{definition}[theo]{Definition}
\newtheorem{theorem}[theo]{Theorem}
\newtheorem*{remas*}{Remarks}
\newtheorem{remas}[theo]{Remarks}
\newtheorem*{exam*}{Example}
\newtheorem{ques}[theo]{Question}
\providecommand{\og}{«\nobreakspace}
\providecommand{\fg}{\nobreakspace»}
\newtheorem{lem}[theo]{Lemma}
\newtheorem{lemma}[theo]{Lemma}
\newtheorem{prop}[theo]{Proposition}
\newtheorem{coro}[theo]{Corollary}
\newtheorem{corol}[theo]{Corollary}
\newtheorem*{conj*}{Conjecture}
\newtheorem{conj}[theo]{Conjecture}
\newtheorem{conjecture}[theo]{Conjecture}
\newtheorem{Proposition}[theo]{Proposition}
\newtheorem{theoreme}[theo]{Theorem}
\theoremstyle{definition}
\newtheorem{defi}[theo]{Definition}
\newtheorem{exam}[theo]{Example}
\newtheorem{exem}[theo]{Example}
\newtheorem{remark}[theo]{Remark}
\newtheorem{example}[theo]{Example}
\newtheorem{rema}[theo]{Remark}
\newtheorem{nota}[theo]{Notation}
\newtheorem{notation}[theo]{Notation}
\newtheorem*{notas}{Notations}
\newtheorem{result}[theo]{Result}
\newtheorem{question}[theo]{Question}
\newtheorem{prob}[theo]{Problem}
\newtheorem*{rema*}{Remark}
\newtheorem*{defi*}{Definition}
\newtheorem*{theo*}{Theorem}
\newtheorem*{lemm*}{Lemma}
\newtheorem*{prop*}{Proposition}
\newcommand{\enoncename}{}
\newtheorem{corpusenonce}[theo]{\enoncename}
\newenvironment{enonce}[2][]{\renewcommand{\enoncename}{#2}\begin{corpusenonce}}{\end{corpusenonce}}
\newenvironment{enonce*}[2][]{\par\medskip\noindent\textbf{#2.}\itshape}{\par\medskip}
\newenvironment{proofc}[1][\proofname]{\begin{proof}[#1]}{\end{proof}}
\newcommand{\Mc}{\mathcal{M}}
\newcommand{\pointir}{.\ ---\ }
\newcommand{\equalenv}[2]{\expandafter\let\csname#1\expandafter\endcsname\csname#2\endcsname\expandafter\let\csname end#1\expandafter\endcsname\csname end#2\endcsname}
\providecommand{\diagbox}[3][]{\begin{tikzpicture}[x=1em,y=1em,baseline=(current bounding box.center)]\useasboundingbox(0,0) rectangle(3,1.4);\draw(0,1.4)--(3,0);\node at(.5,.3){#2};\node at(2.5,1.1){#3};\end{tikzpicture}}
\providecommand{\eg}{e.g.\ }
\providecommand{\ie}{i.e.\ }
\providecommand{\cf}{cf.\ }
\providecommand{\qbinom}[2]{\genfrac{[}{]}{0pt}{}{#1}{#2}}

\providecommand{\mathof}[1]{\mathrm{#1}}

\numberwithin{equation}{section}


\def\endappendix{}



\usepackage{tikz}





\usetikzlibrary{calc}
\usetikzlibrary{decorations.markings}
\usetikzlibrary{decorations.pathreplacing}
\usetikzlibrary{arrows,shapes,positioning}
\tikzstyle directed=[postaction={decorate,decoration={markings,
    mark=at position #1 with {\arrow{>}}}}]
\tikzstyle rdirected=[postaction={decorate,decoration={markings,
    mark=at position #1 with {\arrow{<}}}}]
\tikzset{anchorbase/.style={baseline={([yshift=-0.5ex]current bounding box.center)}}}

\tikzset{
    partial ellipse/.style args={#1:#2:#3}{
        insert path={+ (#1:#3) arc (#1:#2:#3)}
    }
}


\colorlet{green}{black!30!green} 
\newcommand{\END}{{\rm END}}
\newcommand{\Hom}{{\rm Hom}}
\newcommand{\HOM}{{\rm HOM}}
\newcommand{\Gr}{\cat{Flag}_{N}}
\newcommand{\Grn}[1]{\cat{Flag}_{#1}}
\newcommand{\catRep}{{\mathsf{Rep}}}
\newcommand{\SL}{\mathrm{SL}}
\newcommand{\sln}{\mf{sl}_n}
\newcommand{\slm}{\mf{sl}_m}
\newcommand{\slnn}[1]{\mf{sl}_{#1}}
\newcommand{\gln}{\mf{gl}_n}
\newcommand{\glm}{\mf{gl}_m}
\newcommand{\glnn}[1]{\mf{gl}_{#1}}
\def\mf{\mathfrak}
\def\C{{\mathbb C}}
\def\N{{\mathbb N}}
\def\R{{\mathbb R}}
\def\Z{{\mathbb Z}}
\def\Q{{\mathbb Q}}
\def\X{{\mathbb X}}
\newcommand{\rep}{\Rep(\glnn{N})}
\newcommand{\reptwo}{\Rep(\slnn{2})}
\newcommand{\id}{{\rm id}}


\newcommand{\sessAWeb}{N\cat{A}\cat{Web}_\mathrm{s}^{+,\mathrm{ess}}}
\newcommand{\repp}{\Rep^{+}(\glnn{N})}
\newcommand{\reph}{\Rep(\h)}
\newcommand{\rephp}{\Rep^{+}(\h)}
\newcommand{\Sym}{\mathrm{Sym}}\newcommand{\upcev}[1]{\scriptsize\reflectbox{\ensuremath{\vec{\reflectbox{\ensuremath{#1}}}}}}

\def\h{{\mathfrak h}}

\newcommand{\Rep}{\cat{Rep}}
\newcommand{\sh}{w}
\newcommand{\cat}[1]{\ensuremath{\mbox{\bfseries {\upshape {#1}}}}}
\newcommand{\diagrep}{\phi} 
\newcommand{\pTr}{\mathrm{pTr}}
\newcommand{\bVn}{\bigwedge}
\newcommand{\Kar}{\operatorname{Kar}}
\newcommand{\sym}{\mathrm{Sym}}
\newcommand{\Kh}{\mathsf{Kh}}
\newcommand{\KhR}{\mathsf{KhR}}


\newcommand{\TL}{\cat{TL}}
\newcommand{\Web}[1][]{#1\cat{Web}}
\newcommand{\Webp}[1][]{#1\cat{Web}^{+}}
\newcommand{\Webq}[1][]{#1\cat{Web}_{q}}

\newcommand{\wrap}{D}
\newcommand{\wrapi}{D^{-1}}
\newcommand{\A}{\cat{A}}
\newcommand{\ATL}{\cat{A}\cat{TL}}
\newcommand{\AWeb}[1][]{#1\cat{A}\cat{Web}}
\newcommand{\AWebq}[1][]{#1\cat{A}\cat{Web}_{q}}
\newcommand{\AWebp}[1][]{#1\AWeb^+}

\newcommand{\essAWeb}[1][]{#1\AWeb^{\mathrm{ess}}}
\newcommand{\essAWebp}[1][]{#1\AWeb^{\mathrm{ess},+}}

\newcommand{\essbAWeb}[1][]{\overline{#1\AWeb}^{\mathrm{ess}}}
\newcommand{\essbAWebp}[1][]{\overline{#1\AWeb}^{\mathrm{ess},+}}

\newcommand{\Afoam}[1][]{#1\cat{A}\cat{Foam}}
\newcommand{\AFoam}[1][]{#1\cat{A}\cat{Foam}}

\newcommand{\T}{\cat{T}}
\newcommand{\TWebq}[1][]{#1\T\cat{Web}_{q}}
\newcommand{\Tfoam}[1][]{#1\T\cat{Foam}}
\newcommand{\essTfoam}[1][]{#1\Tfoam^{ess}}
\newcommand{\Tlink}{\T\cat{Link}}
\newcommand{\Ttanweb}{\T\cat{TanWeb}}
\newcommand{\Tfoamred}[1][]{#1\T\cat{Foam}^{\mathrm{red}}}

\newcommand{\Su}{\cat{S}}
\newcommand{\SWebq}[1][]{#1\Su\cat{Web}_{q}}
\newcommand{\SWebA}[1][]{#1\Su\cat{Web}_{A}}
\newcommand{\Sfoam}[1][]{#1\Su\cat{Foam}}
\newcommand{\Sfoamred}[1][]{#1\Su\cat{Foam}^{\mathrm{red}}}
\newcommand{\Sfoamor}[1][]{#1\Su\cat{Foam}^{\mathrm or}}
\newcommand{\SCob}{\Su\cat{Cob}}
\newcommand{\Slink}{\Su\cat{Link}}
\newcommand{\Slinko}{\Su\cat{Link}^{\circ}}
\newcommand{\Stanweb}{\Su\cat{TanWeb}}
\newcommand{\Stanwebo}{\Su\cat{TanWeb}^{\circ}}


\newcommand{\repr}{\mathrm{K_0}(\Rep(\glnn{2}))}
\newcommand{\reprs}{\mathrm{K_0}(\Rep(\slnn{2}))}



\newcommand{\cev}[1]{\reflectbox{\ensuremath{\vec{\reflectbox{\ensuremath{#1}}}}}}
\newcommand{\stwo}{*_{\wedge}}
\newcommand{\w}[1]{\wedge^{(#1)}}
\def\la{\langle}
\def\ra{\rangle}















 







  




\begin{document}
\section*{\raggedright 828: Faithfulness of the essential-circle quotient annular glN web functor at q equals one over C}
\noindent\textbf{Author:} Queffelec, Hoel; Wedrich, Paul\par\noindent\textbf{Source:} Extremal weight projectors II, \$\textbackslash{}mathfrak\{gl\}\_\{N\}\$ case. Algebraic Combinatorics 7 (2024), publisher version, DOI 10.5802/alco.330. \url{https://alco.centre-mersenne.org/articles/10.5802/alco.330/}
\par\noindent\textbf{License:} CC BY 4.0, \url{https://creativecommons.org/licenses/by/4.0/}. The original publisher PDF cover has the explicit license notice. See the saved source/license records.
\par\noindent\textbf{Editorial material note (not author text):} The selected problem is Theorem2.35 for the exact q=1 complex essential-circle quotient annular glN web functor. Full native Sections1--2 retain quotient relations, annular reduction, wrap eigenprojectors, their crossing-absorption/local proofs, complete matrix-unit spanning argument, split/merge injectivity and final invertible bending extension. Later categorification/Newton and special gl2 results are unselected and their introduction pointers are retained as original printed locators. Every original web diagram and all mathematical proof text/formulas remain editable native TikZ/TeX; no diagram-category relation is inferred or substituted. Source notation note: the matrix-unit display2139 uses bare sigma immediately after the explicit sigma-epsilon/e-prime choice. The spanning paragraph2149 prints wrap range1,...,n-1 alongside the explicit all-n projector product, while2137 defines all n wraps and2147 uses them in the epimorphism. Both literal discrepancies and the complete surrounding author argument are retained without range/symbol replacement; no editor proof or broader relation is supplied. All proof text and formulas below are editable author TeX. Mathematical wording, language and proof order are retained. Only the article class, title matter and bibliography presentation are adapted for independent compilation; no proof has been generated or repaired.
\par\noindent\textbf{Editorial difficulty:} H3: The proof produces diagrammatic matrix units from wrap-eigenspace projectors, proves their spanning property by annular tangle reduction and same-weight crossing absorption, and extends injectivity through split/merge and bending operations. This explicit diagram-category presentation obligation is research-level; the q=1 complex specialization and essential-circle quotient are essential.
\medskip\hrule\medskip
\noindent\textbf{Author statement, complete proof and local supporting text follow in source order.}
\section{Introduction}



The topological motivation for this article is the search for an extension of
Khovanov--Rozansky link homologies \cite{Kh1, KhR} to invariants of links in
3-manifolds other than $\R^3$. Since quantum link homologies, in their mode of
definition and computation, currently depend on the presentation of links as
2-dimensional projections, the most accessible 3-manifolds in this endeavor are
thickened surfaces $\Sigma\times I$. Just as Khovanov homology categorifies the
Jones polynomial, the surface link homologies should categorify surface skein
modules \cite{Prz1,Tur, APS}. These admit an algebra structure induced by
stacking links, with distinguished bases (conjecturally) satisfying strong
integrality and positivity properties \cite{FG, Thu, Le, MS, Bou}. They provide
quantizations of surface character varieties that play an important role in
quantum Teichm{\"u}ller theory \cite{BW1}. Both aspects make such \textit{skein
algebras} prime targets for categorification via link homology technology.

In order to categorify quantum invariants, it is useful to have explicit,
combinatorial or diagrammatic descriptions of underlying representation
categories. For example, Khovanov homology can be built from a categorification
of the Temperley--Lieb category, which describes the representation category of
$U_q(\slnn{2})$, and all incarnations of Khovanov--Rozansky homology implicitly
employ a categorification of the \textit{MOY} or \textit{web calculus} for the
representation category of $U_q(\glnn{N})$ \cite{MOY, CKM}.

The main purpose of the present paper is to provide representation-theoretic
tools for categorifying skein algebras. The key novelty when working with skein
algebras is that their proposed distinguished bases are obtained from links
colored not by irreducible representations of the corresponding quantum group,
but colored only by the sums of their \textit{extremal weight spaces}, i.e. the
weight spaces for weights in the Weyl group orbit of the highest weight. In
\cite[Section 1.3]{QW} we have proposed a strategy for lifting these colorings
to the categorified level of toric link homologies, which is inspired by
Khovanov's categorification of the colored Jones polynomial~\cite{Kh7}. The main
tool necessary in this approach is a diagrammatic presentation of a certain
subcategory of representations of a Cartan subalgebra $U(\mathfrak{h})\subset
U(\glnn{N})$, which is the first result in this paper:

\begin{result}
In Corollary~\ref{cor:diagrep} we prove that there exists a
diagrammatic presentation for a suitable full subcategory of the representation
category of $U(\mathfrak{h})$, given by an affine extension $\essAWeb[N]$ of the
web calculus for $U(\glnn{N})$ in which $\bVn^k(V)$-labeled essential circles
are set to zero for $0<k<N$.
\end{result}

We emphasize that we deal with universal enveloping algebras, rather than their
quantizations. The former are intended to control the $\C$-linear \textit{morphism
spaces} in our skein module categorifications \cite[Section 1.3]{QW}, whereas
the expected quantum parameter $q$ is promoted to a grading on \textit{objects}.
This is related to the fact that \textit{annular Khovanov homology} \cite{GLW}
has a natural action of $U(\slnn{2})$, not of $U_q(\slnn{2})$. An additional
quantization seems possible, see Remark~\ref{rem:q}, but will not play a role
here. For a detailed discussion of the relation between our diagrammatic
presentation and one that follows from the work of Cautis--Kamnitzer
\cite[Section 2.6]{CK_ann}, we refer to the paper by Lacabanne--Tubbenhauer--Vaz
\cite{LTV} that appeared during revisions on this paper.

In \cite{QW}, we have constructed a diagrammatic presentation for the case of
$\slnn{2}$, and identified idempotent morphisms that encode the projections onto
sums of extremal weight spaces in finite-dimensional
$U(\slnn{2})$-representations. These \textit{extremal weight projectors} are
analogous to, but finer than Jones--Wenzl projectors~\cite{Jon2,Wen} (every
Jones--Wenzl projector splits into a non-negative linear combination of extremal
weight projectors), and they can also be defined recursively. In this article,
we identify and study extremal weight projectors for $\glnn{N}$. 

\begin{result} In Section~\ref{sec:EWP} we give an explicit recursive
  description of the idempotents in the diagrammatic category $\essAWeb[N]$ that
  correspond to projections onto the extremal weight spaces in the
  $U(\glnn{N})$-representations $\sym^k(V)$. We call them extremal weight projectors.
\end{result} 

In fact we prove a slightly stronger version of this result in a central
extension of $\essAWeb[N]$, which has an additional grading by winding number,
that will be important for categorifying skein modules, see
Theorem~\ref{thm:Tm}.

The $\slnn{2}$ extremal weight projectors can be considered as categorifications
of Chebyshev polynomials of the first kind by decategorifying their images to
elements of the representation ring $K_0(\reptwo)\cong \Z[X]$. Analogously, the
extremal weight projectors for $\glnn{N}$ categorify power-sum symmetric
polynomials in the representation ring of $\glnn{N}$. Such categorifications of
classical orthogonal polynomials are of independent interest, see e.g.
\cite{KSa}. Motivated by this, we prove a categorified Newton identity.

\begin{result} In Theorem~3.5 we prove that the extremal weight
projectors satisfy a categorified version of the Newton identity relating
power-sum symmetric and elementary symmetric polynomials.
\end{result} 

The main application for our diagrammatic extremal weight projectors, however,
is in categorifying toric skein modules. In \cite{QW3}, we construct a
categorification of the $\glnn{2}$ skein module of the thickened torus via a
toric $\glnn{2}$ foam category, whose homotopy category is the target of a toric
link homology functor. (The use of $\glnn{2}$ foams, as opposed to Bar-Natan
cobordisms \cite{BN2}, is necessary to guarantee the functoriality of the
resulting link homology, see \cite{Blan, ETW}.) The category $\essAWeb[2]$
describes morphism spaces in this foam category, with affine webs corresponding
to rotationally symmetric foams. In particular, the rotation foams generated by
extremal weight projectors represent indecomposable objects in the Karoubi
completed toric foam category. In Section~4, we prove a delooping
lemma for $\glnn{2}$ webs and decomposition formulas for tensor products of
$\glnn{2}$ extremal weight projectors, which are used in \cite{QW3} to show that
the toric $\glnn{2}$ foams indeed categorify the $\glnn{2}$ skein module of the
thickened torus. We expect that rotation foams generated by the extremal weight
projectors for $\glnn{N}$ give indecomposable objects in a candidate
categorification of the $\glnn{N}$ skein algebra of the torus based on
$\glnn{N}$ foam categories. 



\begin{rema} Affine web categories have appeared before in work of the
first-named author \cite{Queff_aff} on skein modules, and of Cautis--Kamnitzer
\cite{CK_ann} on a K-theoretic version of the derived geometric Satake
correspondence for $SL_N$. The main differences are that here we work at $q=1$,
which makes the affine web categories symmetric monoidal, and that we take a
quotient by $\bVn^k(V)$-labeled essential circles for $0<k<N$. It is unclear to
us how to define an analogous quotient for generic $q$ that would admit extremal
weight projectors.

Affine web categories for $\glnn{N}$ can also be seen as (idempotent completions
of) quotients of HOMFLYPT skein categories as studied, for example in
\cite[Section 4]{Bru}.
\end{rema}
\begin{rema}
\label{rem:q}
Affine web categories at generic $q$ describe morphism spaces in quantized toric
foam categories, which can be defined using a quantized horizontal trace
construction. This is analogous to the quantized annular Bar-Natan cobordisms of
Beliakova--Putyra--Wehrli \cite{BPW}. However, such a quantization involves a
non-canonical choice of a simple closed curve on the torus, that breaks a
natural mapping class group action which is desirable for categorified skein
modules. In \cite{QW3}, we thus proceed with affine webs at $q=1$ and
unquantized toric foam categories.
\end{rema}





\section{Affine \texorpdfstring{$\glnn{N}$}{gl(N)} webs and extremal weight projectors}
We start by recalling the diagrammatic calculus of $\glnn{N}$ webs, which describes the category of representations of $U_q(\glnn{N})$ that is monoidally generated by exterior powers of the vector representation and their duals. 
\subsection{The category of \texorpdfstring{$\glnn{N}$}{gl(N)} webs}
\label{sec:webs}
The category $\Webq[N]$ of $\glnn{N}$ webs is the $\C(q)$-linear pivotal tensor
category with objects generated by points on the line $\R$ that are labeled with
integers in the set $\{1,\dots,N\}$ and that carry an orientation \textit{up} or
\textit{down}. We may consider the tensor unit as a $0$-labeled point without
orientation. The morphisms in $\Webq[N]$ are spanned by \textit{webs} that are
properly embedded in the strip $\R\times [0,1]$, and composition is given by
(the bilinear extension) of stacking in the interval direction. Webs are
trivalent graphs, with edges oriented and labeled by elements of
$\{1,\dots,N\}$, with a flow condition at each vertex imposing that the sum of
incoming labels equals the sum of the outgoing ones. When properly embedded in
the strip, these graphs are interpreted as mapping from the bottom sequence of
boundary points (with labels and orientations) to that at the top. The morphisms
in $\Webq[N]$, represented by $\C(q)$-linear combinations of webs, are
considered modulo isotopy relative to the boundary and subject to certain local
relations. To describe the relations, it is convenient to temporarily allow
edges labeled by $0$ and trivalent vertices involving such edges. The relations
make use of \textit{quantum integers} and \textit{quantum binomial coefficients}
defined for $m\in \N$, $n \in \Z$ as $[n]:=\frac{q^n-q^{-n}}{q-q^{-1}}$,
$[n]!=\prod_{i=1}^n [i]$, and $\qbinom{n}{m}:=\frac{[n]!}{[m]![n-m]!}$. Locally, the following relations hold:

\begin{gather}
\label{eq:webrel}  
\begin{tikzpicture}[anchorbase,scale=.7]
    \draw [very thick,directed=.55] (0,0) to (0,.5); 
    \draw [very thick,directed=.55] (0,1.5) to (0,2); 
   \draw [very thick,directed=.55] (0,.5) to  [out=45,in=315] (0,1.5);
    \draw [very thick,directed=.55] (0,.5) to  [out=135,in=225] (0,1.5);
\node at (.5,1) {\tiny$l$} ;
\node at (-.5,1) {\tiny$k$} ;   
  \end{tikzpicture}
  \;\;=\;\; \qbinom{k+l}{k}\;
  \begin{tikzpicture}[anchorbase,scale=.7]
    \draw [very thick] (0,0) -- (0,2);
    \node at (.45,.25) {\tiny$k+l$} ;    
  \end{tikzpicture}
  \;,\;\;\;
  \begin{tikzpicture}[anchorbase,scale=.7]
    \draw [very thick,directed=.55,directed=.9,directed=.2] (0,0) to (0,2); 
   \draw [very thick,directed=.55] (0,1.25) to  [out=90,in=180] (.25,1.5) to [out=0,in=90] (.5,1.25) to (.5,.75) to  [out=-90,in=0] (.25,.5) to [out=180,in=-90] (0,.75);
\node at (.75,1) {\tiny$l$} ;
\node at (.25,.25) {\tiny$k$} ;   
  \end{tikzpicture}
  \;\;=\;\; \qbinom{N-k}{l}\;
  \begin{tikzpicture}[anchorbase,scale=.7]
    \draw [very thick] (0,0) -- (0,2);
    \node at (.25,.25) {\tiny$k$} ;    
  \end{tikzpicture}
  \;\;,\;\;  
\begin{tikzpicture}[anchorbase,scale=.7]
    \draw [very thick,directed=.55] (0,0) to [out=90,in=225] (0.5,.75);
        \draw [very thick,directed=.55] (1,0) to [out=90,in=315] (0.5,.75);   
    \draw [very thick,directed=.55] (2,0) to [out=90,in=315] (1,1.5);
     \draw [very thick,directed=.55] (0.5,0.75) to [out=90,in=225] (1,1.5);
     \draw [very thick,directed=.55] (1,1.5) to (1,2);
  \end{tikzpicture}
  \;\;=\;\;
  \begin{tikzpicture}[anchorbase,scale=.7]
    \draw [very thick,directed=.55] (0,0) to [out=90,in=225] (1,1.5);
        \draw [very thick,directed=.55] (1,0) to [out=90,in=225] (1.5,.75);   
    \draw [very thick,directed=.55] (2,0) to [out=90,in=315] (1.5,.75);
     \draw [very thick,directed=.55] (1.5,0.75) to [out=90,in=315] (1,1.5);
     \draw [very thick,directed=.55] (1,1.5) to (1,2);
  \end{tikzpicture}
\\
\nonumber
\begin{tikzpicture}[anchorbase,scale=.75]
  \draw [very thick,->] (0,0) -- (0,1.5);
  \draw [very thick, ->] (1,0) -- (1,1.5);
  \draw [very thick, directed=.5] (0,.4) -- (1,.6);
  \draw [very thick, directed=.5] (1,.9) -- (0,1.1);
  \node at (.5,.2) {\tiny $b$};
  \node at (.5,1.2) {\tiny $a$};
  \node at (0,-.2) {\tiny $k$};
  \node at (1,-.2) {\tiny $l$};
\end{tikzpicture} 
\;\;=\;\;
\sum_{t} \qbinom{k-l+a-b}{t}\;\;
\begin{tikzpicture}[anchorbase,scale=.75]
  \draw [very thick,->] (0,0) -- (0,1.5);
  \draw [very thick, ->] (1,0) -- (1,1.5);
  \draw [very thick, directed=.5] (1,.4) -- (0,.6);
  \draw [very thick, directed=.5] (0,.9) -- (1,1.1);
  \node at (.5,.2) {\tiny $a-t$};
  \node at (.5,1.2) {\tiny $b-t$};
  \node at (0,-.2) {\tiny $k$};
  \node at (1,-.2) {\tiny $l$};
\end{tikzpicture}
\quad,\quad
\begin{tikzpicture}[anchorbase,scale=.75]
  \draw [very thick,->] (0,0) -- (0,1.5);
  \draw [very thick, <-] (1,0) -- (1,1.5);
  \draw [very thick, rdirected=.5] (0,.6) to [out=330,in=210](1,.6);
  \draw [very thick, rdirected=.5] (1,.9) to [out=150,in=30]  (0,.9);
  \node at (.5,.1) {\tiny $1$};
  \node at (.5,1.3) {\tiny $1$};
  \node at (0,-.2) {\tiny $k$};
  \node at (1,-.2) {\tiny $l$};
\end{tikzpicture}
\;=\;
\begin{tikzpicture}[anchorbase,scale=.75]
  \draw [very thick,->] (0,0) -- (0,1.5);
  \draw [very thick, <-] (1,0) -- (1,1.5);
  \draw [very thick, rdirected=.5] (1,.4) to [out=150,in=30] (0,.4);
  \draw [very thick, rdirected=.5] (0,1.1) to [out=330,in=210] (1,1.1);
  \node at (.5,.2) {\tiny $1$};
  \node at (.5,1.2) {\tiny $1$};
  \node at (0,-.2) {\tiny $k$};
  \node at (1,-.2) {\tiny $l$};
\end{tikzpicture}
+
[N-k-l] \begin{tikzpicture}[anchorbase,scale=.75]
  \draw [very thick,->] (0,0) -- (0,1.5);
  \draw [very thick, <-] (1,0) -- (1,1.5);
  \node at (0,-.2) {\tiny $k$};
  \node at (1,-.2) {\tiny $l$};
\end{tikzpicture}
\end{gather}
as well as their orientation reversals. The following relations are useful
consequences of the ones above:
\begin{equation}
\label{eq:webrel2}
\begin{tikzpicture}[anchorbase,scale=.75]
  \draw [very thick,directed=.5] (0,0) circle (.5);
  \node at (.6,-.25) {\tiny $l$};
\end{tikzpicture}
\;\; = \;\;
\begin{tikzpicture}[anchorbase,scale=.75]
  \draw [very thick,rdirected=.5] (0,0) circle (.5);
  \node at (.6,-.25) {\tiny $l$};
\end{tikzpicture}
\;\;=\;\;\qbinom{N}{l}\quad, \quad 
\begin{tikzpicture}[anchorbase,scale=.5]
\draw [very thick,->] (0,0) to  (0,2);
\draw [very thick,->] (1,2) to (1,0);
  \node at (0,-.2) {\tiny $N$};
  \node at (1,-.2) {\tiny $N$};
\end{tikzpicture}
\;\; =\;\;
\begin{tikzpicture}[anchorbase,scale=.5]
\draw [very thick,->] (0,0) to (0,.5) to [out=90,in=90] (1,.5) to (1,0);
\draw [very thick,->] (1,2) to (1,1.5) to [out=270,in=270] (0,1.5) to (0,2);
  \node at (0,-.2) {\tiny $N$};
  \node at (1,-.2) {\tiny $N$};
\end{tikzpicture}
\quad , \quad
\begin{tikzpicture}[anchorbase,scale=.5]
\draw [very thick,->] (0,0) to  (0,2);
\draw [very thick,->] (1,2) to (1,0);
  \node at (0,-.2) {\tiny $k$};
  \node at (1,-.2) {\tiny $N$};
\end{tikzpicture}
\;\; = \;\;
\begin{tikzpicture}[anchorbase,scale=.5]
\draw [very thick,->] (0,0) to (0,0.25) to [out=90,in=240] (0.5,0.75) to [out=120,in=240] (0.5,1.25) to[out=120,in=270] (0,1.75) to  (0,2);
\draw [very thick,<-] (1,0) to (1,0.25) to [out=90,in=0] (.5,.75) ;
\draw [very thick] (1,2) to (1,1.75) to [out=270,in=0] (0.5,1.25);
  \node at (0,-.2) {\tiny $k$};
  \node at (1,-.2) {\tiny $N$};
   \node at (1.25,1) {\tiny $N-k$};
\end{tikzpicture}
\end{equation}

The first equation is a special case of the second relation in
\eqref{eq:webrel}, obtained by setting $k=0$ and erasing $0$-labeled edges.








\subsection{Link with \texorpdfstring{$U_q(\glnn{N})$}{Uq(gl(N)}-representation theory}
% ---------------------------------------------------------------------------------------
Let $\Rep(U_q(\glnn{N}))$ denote the $\C(q)$-linear pivotal tensor category of
$U_q(\glnn{N})$-representations that is monoidally generated by quantum exterior powers
of the vector representation and their duals. The main purpose of the
diagrammatic calculus of $\glnn{N}$ webs is to describe this category.

\begin{theo}\label{thm:CKM} There exists an equivalence of $\C(q)$-linear pivotal tensor categories
\[\phi \colon \Webq[N] \to \Rep(U_q(\glnn{N}))\]
that sends $k$-labeled upward points to $k$-fold exterior powers of the vector representation of $U_q(\glnn{N})$.
\end{theo}
Essentially, this theorem is due to Cautis--Kamnitzer--Morrison~\cite{CKM},
although they state it for $U_q(\slnn{N})$ in \cite[Theorem 3.3.1]{CKM}.
Versions for $U_q(\glnn{N})$ have appeared as special cases of \cite[Main Theorem A]{QS2} and \cite[Theorem 3.20]{TVW}. We now describe the functor $\phi$
explicitly.

Recall that $U_q(\glnn{N})$ is the $\C(q)$-algebra generated by $E_i,F_i$ for $1\leq i \leq N-1$ and $L^{\pm 1}_j$ for $1\leq j \leq N$ subject to the following relations:
\begin{gather}
L_i E_i=q E_i L_i,\quad L_iF_i=q^{-1}F_iL_i,\quad L_{i+1} E_i=q^{-1} E_iL_{i+1},\quad L_{i+1}F_i=q F_i L_{i+1}
\\
[E_i,F_j]=\delta_{i,j} \frac{L_iL_{i+1}^{-1}-L_{i+1}L_{i}^{-1}}{q-q^{-1}},\quad [L_i,L_j]=0.
\\
E_i^2 E_j -[2] E_i E_j E_i + E_j E_i^2 =0 \text{ if } |i-j|=1  \text{ and } [E_i,E_j]=0 \text{ otherwise;}
\\
\text{analogously for $F_i$s}. \nonumber
\end{gather}
It is a Hopf algebra with coproduct, antipode and counit as follows:
\begin{gather*}
  \Delta(E_i)=E_i\otimes L_iL_{i+1}^{-1} +1\otimes E_i, \quad \Delta(F_i)=F_i\otimes 1 +L_i^{-1}L_{i+1}\otimes F_i \\
  \Delta(L_i^{\pm 1})=L_i^{\pm 1}\otimes L_i^{\pm 1}\\
S(L_i^{\pm 1})=L_i^{\mp 1}, \quad S(E_i)=-E_i L_i^{-1}L_{i+1}, \quad S(F_i)= - L_iL_{i+1}^{-1}F_i \\
\epsilon(L_i^{\pm 1})= 1, \quad \epsilon(E_i)=0, \quad \epsilon(F_i)=0
\end{gather*}
Let $V=\C(q)\langle v_1,v_2,\dots v_N\rangle$ denote the vector representation of $U_q(\glnn{N})$ and $V^*=\C(q)\langle v_1^*,v_2^*,\dots,v_N^* \rangle$ its dual. 

To leftward oriented cups and caps, the functor $\phi$ associates the natural evaluation and co-evaluation maps for duals:
\begin{align*}
\begin{tikzpicture}[anchorbase, scale=.5]
\draw [very thick,->] (1,2) to (1,1.7)to [out=270,in=0] (.5,1.2) to [out=180,in=270] (0,1.7) to (0,2);
\end{tikzpicture}
\quad \xmapsto{\phi} \quad
\begin{cases}
\C \to  V \otimes V^* \\
1\mapsto  \sum_{k=1}^N v_k\otimes v_k^*
\end{cases}
\quad ,\quad
\begin{tikzpicture}[anchorbase, scale=.5]
\draw [very thick,->] (1,0) to (1,0.3)to [out=90,in=0] (.5,.8) to [out=180,in=90] (0,.3) to (0,0);
\end{tikzpicture}
\quad \xmapsto{\phi} \quad
\begin{cases}
V^*\otimes V \to \C \\
v_{k}^*\otimes v_{l} \mapsto \delta_{k,l}
\end{cases}
\end{align*}
Let $\bVn^k V$ denote the $k$-th quantum exterior power of $V$. This has a basis
indexed by subsets $S=\{i_1,\dots,i_k \}\subset \{1,2,\dots, N\}$ of size $k$.
If $1\leq i_1<\cdots < i_k\leq N$, we use the following notation for the
corresponding basis vector: $v_S:=v_{i_1}\wedge \cdots \wedge v_{i_k}$. The dual
$(\bVn^k V)^*\cong \bVn^k V^*$ then has the dual basis given by vectors $v_S^*$
and we have corresponding thickness $k$ cap and cup morphisms as above. 

\begin{align*}
\begin{tikzpicture}[anchorbase, scale=.5]
\draw [very thick,->] (1,2) to (1,1.7)to [out=270,in=0] (.5,1.2) to [out=180,in=270] (0,1.7) to (0,2);
\node at (.5,1.5) {\tiny$k$} ;
\end{tikzpicture}
\quad \xmapsto{\phi} \quad
\begin{cases}
\C \to  \bVn^k V \otimes \bVn^k V^* \\
1\mapsto  \sum_{|S|=k} v_S\otimes v_S^*
\end{cases}
\quad ,\quad
\begin{tikzpicture}[anchorbase, scale=.5]
\draw [very thick,->] (1,0) to (1,0.3)to [out=90,in=0] (.5,.8) to [out=180,in=90] (0,.3) to (0,0);
\node at (.5,.5) {\tiny$k$} ;
\end{tikzpicture}
\quad \xmapsto{\phi} \quad
\begin{cases}
\bVn^k V^*\otimes \bVn^k V \to \C \\
v_{S}^*\otimes v_{T} \mapsto \delta_{S,T}
\end{cases}
\end{align*}

The other duality maps, that is, rightward oriented caps and cups, are perturbed by powers of $q$. 
\begin{align*}
\begin{tikzpicture}[anchorbase, scale=.5]
\draw [very thick,<-] (1,2) to (1,1.7)to [out=270,in=0] (.5,1.2) to [out=180,in=270] (0,1.7) to (0,2);
\node at (.5,1.5) {\tiny$k$} ;
\end{tikzpicture}
\quad \xmapsto{\phi} \quad
\begin{cases}
\C \to  \bVn^k V^* \otimes \bVn^k V \\
1\mapsto  \sum_{|S|=k} q^{-\epsilon_S}v_S^*\otimes v_S
\end{cases}
\quad ,\quad
\begin{tikzpicture}[anchorbase, scale=.5]
\draw [very thick,<-] (1,0) to (1,0.3)to [out=90,in=0] (.5,.8) to [out=180,in=90] (0,.3) to (0,0);
\node at (.5,.5) {\tiny$k$} ;
\end{tikzpicture}
\quad \xmapsto{\phi} \quad
\begin{cases}
\bVn^k V\otimes \bVn^k V^* \to \C \\
v_{S}\otimes v_{T}^* \mapsto \delta_{S,T} q^{\epsilon_S}
\end{cases}
\end{align*}
Here $\epsilon_S= \sum_{i\in S}(N+1-2i)$.


Merges of thick strands act as ($q$-deformed) exterior product:
\begin{align} \label{eq:thickmerge}
\begin{tikzpicture}[anchorbase, scale=.5]
\draw [very thick, directed=.55] (0,0) to [out=90,in=225] (.5,.7);
\draw [very thick, directed=.55] (1,0) to [out=90,in=315]  (.5,.7);
\draw [very thick,->] (.5,.7) -- (.5,1.4);
\node at (-.2,-.2) {\tiny$k$} ;
\node at (1.2,-.2) {\tiny$l$} ;
\node at (.5,1.6) {\tiny$k+l$} ;
\end{tikzpicture}
\quad \xmapsto{\phi} \quad
\begin{cases}
\bVn^k V\otimes \bVn^l V \to \bVn^{k+l} V \\
v_{S}\otimes v_{T} \mapsto 0 & \text{ if } S\cap T\neq \emptyset \\
v_{S}\otimes v_{T} \mapsto (-q)^{\epsilon_{S,T}} v_{S\cup T} & \text{ otherwise } 
\end{cases}
\end{align}
Here $\epsilon_{S,T}$ is the number of inversions in the concatenation of the ordered lists of elements of $S$ and $T$. The split vertex acts as follows:
\begin{align} \label{eq:thicksplit}
\begin{tikzpicture}[anchorbase,scale=.5]
\draw [very thick, directed=.5] (.5,0) -- (.5,.7);
\draw [very thick, directed=1] (.5,.7) to [out=135,in=270]  (0,1.3)to (0,1.4);
\draw [very thick, directed=1] (.5,.7) to [out=45,in=270]  (1,1.3) to (1,1.4);
\node at (-.2,1.6) {\tiny$k$} ;
\node at (1.2,1.6) {\tiny$l$} ;
\node at (.5,-.2) {\tiny$k+l$} ;
\end{tikzpicture}
\quad \xmapsto{\phi}\quad
\begin{cases}
\bVn^{k+l} V  \to \bVn^{k} V \otimes \bVn^{l} V  \\
v_S \mapsto (-1)^{k l}\sum_{T\subset S, |T|=k} (-q)^{-\epsilon_{S\setminus T,T}} v_T\otimes v_{S\setminus T}
\end{cases}
\end{align}

Analogous formulas hold for merges and splits of duals, which implies that merges and splits can be slid around caps and cups:
\begin{align}
\begin{tikzpicture}[anchorbase, scale=.5]
\draw [very thick, <-] (0,0) to [out=90,in=225] (.5,.7);
\draw [very thick, <-] (1,0) to [out=90,in=315]  (.5,.7);
\draw [very thick,rdirected=.55] (.5,.7) -- (.5,1.4);
\node at (-.2,-.2) {\tiny$k$} ;
\node at (1.2,-.2) {\tiny$l$} ;
\node at (.5,1.6) {\tiny$k+l$} ;
\end{tikzpicture}
\quad \xmapsto{\phi} \quad
\begin{cases}
\bVn^k V^*\otimes \bVn^l V^* \to \bVn^{k+l} V^* \\
v_{S}^*\otimes v_{T}^* \mapsto 0 & \text{ if } S\cap T\neq \emptyset \\
v_{S}^*\otimes v_{T}^* \mapsto (-1)^{k l}(-q)^{-\epsilon_{S,T}} v_{S\cup T}^* & \text{ otherwise } 
\end{cases}
\end{align}
and:
\begin{align}
\begin{tikzpicture}[anchorbase,scale=.5]
\draw [very thick, <-] (.5,0) -- (.5,.7);
\draw [very thick, rdirected=.55] (.5,.7) to [out=135,in=270]  (0,1.3)to (0,1.4);
\draw [very thick, rdirected=.55] (.5,.7) to [out=45,in=270]  (1,1.3) to (1,1.4);
\node at (-.2,1.6) {\tiny$k$} ;
\node at (1.2,1.6) {\tiny$l$} ;
\node at (.5,-.2) {\tiny$k+l$} ;
\end{tikzpicture}
\quad \xmapsto{\phi}\quad
\begin{cases}
\bVn^{k+l} V^*  \to \bVn^{k} V^* \otimes \bVn^{l} V^*  \\
v_S^* \mapsto \sum_{T\subset S, |T|=k} (-q)^{\epsilon_{S\setminus T,T}} v_T^*\otimes v_{S\setminus T}^*
\end{cases}
\end{align}
This concludes the description of the functor $\phi$.

The category $\Rep(U_q(\glnn{N}))$ is braided, and by virtue of Theorem~\ref{thm:CKM}, so is $\Webq[N]$. The diagrammatic description of the braiding of two fundamental $U_q(\glnn{N})$-representations in $\Webq[N]$ is given as follows:
\begin{equation} \label{eq:crossing}
\begin{tikzpicture}[anchorbase,scale=.75]
  \draw [very thick, ->] (1,0) -- (0,1.5);
  \draw [white, line width=.15cm] (0,0) -- (1,1.5);
  \draw [very thick, ->] (0,0) -- (1,1.5);
  \node at (0,-.3) {\tiny $k$};
  \node at (1,-.3) {\tiny $l$};
\end{tikzpicture}
\quad = \quad
(-q)^{kl}\sum_{\substack{a,b\geq 0\\ b-a=k-l}} (-q)^{b-k}\quad
\begin{tikzpicture}[anchorbase,scale=.75]
  \draw [very thick,->] (0,0) -- (0,1.5);
  \draw [very thick, ->] (1,0) -- (1,1.5);
  \draw [very thick, directed=.5] (0,.4) -- (1,.6);
  \draw [very thick, directed=.5] (1,.9) -- (0,1.1);
  \node at (.5,.2) {\tiny $b$};
  \node at (.5,1.2) {\tiny $a$};
  \node at (0,-.2) {\tiny $k$};
  \node at (1,-.2) {\tiny $l$};
\end{tikzpicture}
\end{equation}


In particular, a crossing of two $1$-labeled strands is given by:
\[
\begin{tikzpicture}[anchorbase,scale=.75]
  \draw [very thick, ->] (1,0) -- (0,1.5);
  \draw [white, line width=.15cm] (0,0) -- (1,1.5);
  \draw [very thick, ->] (0,0) -- (1,1.5);
\end{tikzpicture}
\quad = \quad
\begin{tikzpicture}[anchorbase,scale=.75]
  \draw [very thick, ->] (0,0) -- (0,1.5);
  \draw [very thick, ->] (1,0) -- (1,1.5);
\end{tikzpicture}
\quad - \;\;q\;\;
\begin{tikzpicture}[anchorbase,scale=.75]
  \draw [very thick] (0,0) to [out=90,in=-90] (.5,.7);
  \draw [very thick] (1,0) to [out=90,in=-90] (.5,.7);
  \draw [very thick, directed=1] (.5,.7) -- (.5,.8);
  \node at (.75,.75) {\tiny $2$};
  \draw [very thick,->] (.5,.8) to [out=90,in=-90] (0,1.5);
  \draw [very thick,->] (.5,.8) to [out=90,in=-90] (1,1.5);
\end{tikzpicture}
\]
For negative crossings, one uses the above formulas with $q$ inverted. (Setting
$q=1$ recovers the ordinary braiding given by swapping tensor factors.)

\begin{lemma} \label{lem:Reid} The following analogs of Reidemeister moves hold in $\Webq[N]$, where strands can carry all possible orientations and labels.
\begin{gather*}
  q^{-k(N-1)} \begin{tikzpicture}[anchorbase,scale=.7]
   \draw [very thick] (0,0) to [out=90,in=180] (.75,1.25) to [out=0,in=90] (1.25,.75); 
     \draw [white, line width=.15cm]  (1.25,.75) to [out=-90,in=0] (.75,.25) to [out=180,in=-90] (0,1.5);     
      \draw [very thick]  (1.25,.75) to [out=-90,in=0] (.75,.25) to [out=180,in=-90] (0,1.5);     
        \node at (0,-.3) {\tiny $k$};
  \end{tikzpicture}
  \;\;=\;\;
  \begin{tikzpicture}[anchorbase,scale=.7]
    \draw [very thick] (0,0) -- (0,1.5); 
     \node at (0,-.3) {\tiny $k$};
  \end{tikzpicture}
  \;\;=\;\;
  q^{k(N-1)} 
    \begin{tikzpicture}[anchorbase,scale=.7]
      \draw [very thick]  (1.25,.75) to [out=-90,in=0] (.75,.25) to [out=180,in=-90] (0,1.5);     
       \draw [white,line width=.15cm] (0,0) to [out=90,in=180] (.75,1.25) to [out=0,in=90] (1.25,.75);   
        \draw [very thick] (0,0) to [out=90,in=180] (.75,1.25) to [out=0,in=90] (1.25,.75);   
         \node at (0,-.3) {\tiny $k$};
  \end{tikzpicture}, 
  \\
  \begin{tikzpicture}[anchorbase,scale=.5]
    \draw [very thick] (0,0) to [out=90,in=-90] (1,1) to [out=90,in=-90] (0,2);
    \draw [white, line width=.15cm] (1,0) to [out=90,in=-90] (0,1) to [out=90,in=-90] (1,2);
    \draw [very thick] (1,0) to [out=90,in=-90] (0,1) to [out=90,in=-90] (1,2);
  \end{tikzpicture}
\quad=\quad
  \begin{tikzpicture}[anchorbase,scale=.5]
    \draw [very thick] (0,0) -- (0,2);
    \draw [very thick] (1,0) -- (1,2);
  \end{tikzpicture}
  \quad,\quad
  \begin{tikzpicture}[anchorbase,scale=.5]
    \draw [very thick] (0,0) to [out=90,in=-90] (.7,.7) to [out=90,in=-90] (1.4,1.4) -- (1.4,2.1);
        \draw [white, line width=.15cm] (.7,0) to [out=90,in=-90] (0,.7) -- (0,1.4) to [out=90,in=-90] (.7,2.1);
    \draw [very thick] (.7,0) to [out=90,in=-90] (0,.7) -- (0,1.4) to [out=90,in=-90] (.7,2.1);
        \draw [white, line width=.15cm] (1.4,0) -- (1.4,.7) to [out=90,in=-90] (.7,1.4) to [out=90,in=-90] (0,2.1);
    \draw [very thick] (1.4,0) -- (1.4,.7) to [out=90,in=-90] (.7,1.4) to [out=90,in=-90] (0,2.1);
  \end{tikzpicture}
\quad=\quad
  \begin{tikzpicture}[anchorbase,scale=.5]
    \draw [very thick] (0,0) -- (0,.7) to [out=90,in=-90] (.7,1.4) to [out=90,in=-90] (1.4,2.1);    
    \draw [white, line width=.15cm] (.7,0) to [out=90,in=-90] (1.4,.7) -- (1.4,1.4) to [out=90,in=-90] (.7,2.1);
    \draw [very thick] (.7,0) to [out=90,in=-90] (1.4,.7) -- (1.4,1.4) to [out=90,in=-90] (.7,2.1);
        \draw [white, line width=.15cm] (1.4,0) to [out=90,in=-90] (.7,.7) to [out=90,in=-90] (0,1.4) -- (0,2.1);
            \draw [very thick] (1.4,0) to [out=90,in=-90] (.7,.7) to [out=90,in=-90] (0,1.4) -- (0,2.1);
  \end{tikzpicture}
  \quad,\quad
    \begin{tikzpicture}[anchorbase,scale=.5]
    \draw [very thick] (.5,0) to  (.5,.7) to [out=45,in=-90] (1,1.4) -- (1,2.1);
       \draw [very thick] (.5,.7) to [out=135,in=-90] (0,1.4) -- (0,2.1);
       \draw [white, line width=.15cm] (-.4,1.4) to  (1.4,1.4);
    \draw [very thick] (-.4,1.4) to  (1.4,1.4);
  \end{tikzpicture}
  \quad =\quad
  \begin{tikzpicture}[anchorbase,scale=.5]
    \draw [very thick] (.5,0) to  (.5,1.4) to [out=45,in=-90] (1,2.1);
       \draw [very thick] (.5,1.4) to [out=135,in=-90] (0,2.1);
    \draw [white, line width=.15cm] (-.4,0.7) to  (1.4,0.7);
    \draw [very thick] (-.4,0.7) to  (1.4,0.7);
  \end{tikzpicture}
  \end{gather*}
  We will refer to the last relation as a \textit{forkslide move}.
\end{lemma}
\begin{proof} The Reidemeister II, III and forkslide moves follow from the property of a braiding, and our braiding convention is only a minor rescaling of the one in \cite[Corollary 6.2.3]{CKM}, see also \cite[Section 2.4]{TVW}. The Reidemeister I moves can be verified inductively as in \cite[Lemma 2.9]{Queff_aff}.
\end{proof}
  
\begin{defi}
We denote by $\Webp[N]$ the full subcategory of $\Web[N]$ with objects given by arbitrary sequences with exclusively upward pointing orientations. 
\end{defi}
In the following we will use the same superscript $+$ to indicate analogous full
subcategories of other web categories, consisting of those objects with upward
(or outward) pointing orientations. The next lemma is a well-known consequence
the proof of Theorem~\ref{thm:CKM} using quantum skew Howe duality, see
\cite[Section 4.3 (fullness) and Section 5 (ladder webs)]{CKM} and analogously in \cite[Prop 6.8]{QS2} and \cite{TVW}.

\begin{lemma}\label{lem:upward} The morphism spaces of $\Webp[N]$ are spanned by
upward-pointing webs, i.e. webs whose edges admit oriented parametrisations with
derivative having a positive vertical component everywhere.
\end{lemma} 
  
In the following, we will consider \emph{skein modules} of isotopy classes of
webs as in Section~\ref{sec:webs}, but embedded in different surfaces, modulo
the local relations from \eqref{eq:webrel}, and we will also vary the ground
ring. In the following sections we deal with webs over $\C$, whose defining
relations are obtained from \eqref{eq:webrel} by specializing $q=1$. We indicate
categories of webs at $q=1$ by the omission of the $q$-subscript, e.g. $\Web[N]$
instead of $\Webq[N]$. The functor $\phi$ also specializes to $q=1$ and then
relates $\Web[N]$ to the symmetric monoidal category of
$U(\glnn{N})$-representations.

Note that setting $q=1$ identifies the evaluation of positive and negative crossings in terms of webs in \eqref{eq:crossing}, and so we sometimes do not display any over- or under-crossing information in graphics. In particular, the braid group action induced by 1-labeled crossings becomes a symmetric group action.

  
\subsection{Affinization at \texorpdfstring{$q=1$}{q=1}}
In \cite{QW}, we considered an affine extension of the Temperley-Lieb category, and extended an analog of the functor $\phi$ to this more general category. Just as in this simpler $\slnn{2}$ case, we will consider a more general affine web category that will give us the freedom to extend the diagrammatic presentation of the representation category of $U(\glnn{N})$ to a Cartan subalgebra.


We define the category $\AWeb[N]$ to be the $\C$-linear category with morphisms
spaces spanned by webs properly embedded in the annulus $[0,1]\times
[0,1]/\{(0,s)\sim (1,s)\mid s\in [0,1]\}$, subject to the same local relations
as in $\Web[N]$, i.e. relations \eqref{eq:webrel} at $q=1$. We remember the
\textit{seam} $\alpha=\{(0,s)\sim (1,s)\mid s\in [0,1]\}$, which will be drawn
as a dashed line in illustrations. The endpoints of a web are on the two
boundary circles $[0,1]\times \{0\}\setminus\{(0,0)\}$ and $[0,1]\times
\{1\}\setminus\{(0,1)\} $, so that the web can be interpreted as a mapping
from the configuration of points on the first (inner) circle to the second (outer) one. We further require of webs that their
endpoints are disjoint from the endpoints of the seam.

Just as before, one can compose annular webs by stacking the annuli. It is easy to see that the morphisms of $\AWeb[N]$ can be generated from those
morphisms that are supported in the strip $(0,1)\times [0,1 ]$ (which may be
considered as specifying morphisms in $\Web[N]$), and the additional new
\textit{wrapping morphisms}: 

\[
\wrap 
\;\;=\;\;
\begin{tikzpicture}[anchorbase, scale=.3]
%% Inner boundary circle
\draw (0,0) circle (1);
%% Outer boundary circle
\draw (0,0) circle (3);
%% wrapping around
\draw [very thick] (-.8,.6) .. controls (-2.6,1.2) and (-2,-2) .. (0,-2) .. controls (2,-2) and (1.2,.9) ..  (2.4,1.8);
\draw [very thick] (-.4,.9) .. controls (-.8,1.8) and (-.9,1.2) .. (-1.8,2.4);
\node at (0,2) {$\dots$};
\draw [very thick] (.8,.6) .. controls (1.2,.9) and (.75,1.3) .. (1.5,2.6);
% boundary markings
\node at (0,-1) {$*$};
\node at (0,-3) {$*$};
\draw [dashed] (0,-1) to (0,-3);
\end{tikzpicture}
\quad,\quad
\wrapi
\;\;=\;\;
\begin{tikzpicture}[anchorbase, scale=.3]
%% Inner boundary circle
\draw (0,0) circle (1);
%% Outer boundary circle
\draw (0,0) circle (3);
%% wrapping around
\draw [very thick] (.8,.6) .. controls (2.6,1.2) and (2,-2) .. (0,-2) .. controls (-2,-2) and (-1.2,.9) ..  (-2.4,1.8);
\draw [very thick] (.4,.9) .. controls (.8,1.8) and (.9,1.2) .. (1.8,2.4);
\node at (0,2) {$\dots$};
\draw [very thick] (-.8,.6) .. controls (-1.2,.9) and (-.75,1.3) .. (-1.5,2.6);
% boundary markings
\node at (0,-1) {$*$};
\node at (0,-3) {$*$};
\draw [dashed] (0,-1) to (0,-3);
\end{tikzpicture}
\]


When starting from an abstract annulus rather than our concrete model, a choice
of seam endpoints is necessary to make the composition in $\AWeb[N]$ well-defined
and a choice of seam is necessary to distinguish identity morphisms from all
other (invertible) wrapping morphisms.

We also stress that in $\AWeb[N]$ webs can
come with any orientation on the boundary. 

\subsection{Link with representation theory}
We will extend the domain of the functor $\phi$ from $\Web[N]$ to $\AWeb[N]$ by
sending the wrapping morphisms to maps between $U(\glnn{N})$-representations,
which respect the weight space decomposition but break the $U(\glnn{N})$-action.
This will allow us to build new diagrammatic projectors, and we will now explain
how to choose this preferred extension. 

We first consider a single counterclockwise wrap morphism $\wrap=\wrap_1$ of a single 1-labeled outward pointing strand. %\vspace{-3mm}
\[
%\wrap 
%\;\;=\;\;
\begin{tikzpicture}[anchorbase, scale=.3]
%% Inner boundary circle
\draw (0,0) circle (1);
%% Outer boundary circle
\draw (0,0) circle (3);
%% wrapping around
\draw [very thick,->] (0,1) .. controls (0,1.7) and (-1.6,1) .. (-1.6,0) to [out=-90,in=180] (0,-1.8) to [out=0,in=-90] (2.1,0) to [out=90,in=-90] (0,3);
% boundary markings
\node at (-.2,2) {\tiny$1$};
\node at (0,-3) {$*$};
\draw [dashed] (0,-1) to (0,-3);
\end{tikzpicture}
\]
The requirement that $\phi(\wrap)$ respects the weight space decomposition of
$V$ implies that $\phi(\wrap)(v_k)= \gamma_k v_k$ for some $\gamma_k\in \C$ and
the desired invariance under ambient isotopy forces these scalars to be
invertible. In fact, this choice of scalars determines the action of
$\phi(\wrap_k)$, the $k$-labeled version of the wrap. To see this, note that the
first relation in \eqref{eq:webrel} allows one to open a digon in the $k$-edge.
Iterating this procedure, one can open a $k$-blister in the $k$-edge, i.e. a
configuration of nested digons with $k$-many parallel $1$-edges in the innermost
part (see the proof of Lemma~\ref{lem:resolveastangles} below for an
illustration). Sliding one half of the $k$-blister around the wrap $\wrap_k$,
the eigenvalues of $\phi(\wrap_k)$ can be seen to be $k$-fold products of the
eigenvalues of $\phi(\wrap)$: $\phi(\wrap_k)(v_S)= (\prod_{i\in S}\gamma_i)
v_S$. Furthermore, inverse wraps have inverse eigenvalues:
$\phi(\wrap_k^{-1})(v_S)= (\prod_{i\in S}\gamma_i)^{-1} v_S$. Next, we would
like to have relations of the form:%\vspace{-3mm}

\begin{equation}
\label{eq:capslide}
\phi
\left(
\begin{tikzpicture}[anchorbase, scale=.3]
%% Inner boundary circle
\draw (0,0) circle (1);
%% Outer boundary circle
\draw (0,0) circle (3);
%% wrapping around
\draw [very thick] (-.8,.6) to [out=150,in=90] (-1.75,0) to[out=270,in=180] (0,-2) to [out=0,in=270] (2,0) to [out=90,in
=225] (2.25,.9) to [out=45,in=315] (2.25,1.35)to[out=135,in=45] (1.8,1.35) to (.8,.6) ;
\draw [very thick] (-.4,.9)to  (-1.2,2.7);
\node at (0,2) {$\dots$};
\draw [very thick] (.4,.9)  to (1.2,2.7);
% boundary markings
\node at (0,-1) {$*$};
\node at (0,-3) {$*$};
\draw [dashed] (0,-1) to (0,-3);
\end{tikzpicture}
\right)
\;=\;
\phi\left(
\begin{tikzpicture}[anchorbase, scale=.3]
%% Inner boundary circle
\draw (0,0) circle (1);
%% Outer boundary circle
\draw (0,0) circle (3);
%% wrapping around
\draw [very thick] (.8,.6) to [out=30,in=90] (1.75,0) to[out=270,in=0] (0,-2) to [out=180,in=270] (-2,0) to [out=90,in
=315] (-2.25,.9) to [out=135,in=225] (-2.25,1.35)to[out=45,in=135] (-1.8,1.35) to (-.8,.6) ;
\draw [very thick] (-.4,.9)to  (-1.2,2.7);
\node at (0,2) {$\dots$};
\draw [very thick] (.4,.9)  to (1.2,2.7);
% boundary markings
\node at (0,-1) {$*$};
\node at (0,-3) {$*$};
\draw [dashed] (0,-1) to (0,-3);
\end{tikzpicture}
\right)
\end{equation}
To ensure that $\phi$ respects such isotopy relations for sliding cups and caps around the annulus, we need to have
$\phi(\wrap_k)(v_S^*)= (\prod_{i\in S}\gamma_i)^{-1} v_S^*$ and $\phi(\wrap_k^{-1})(v_S^*)= (\prod_{i\in S}\gamma_i) v_S^*$, which determine the maps assigned to inward pointing versions of $D$ and $D^{-1}$. 



In order to be able to project onto the 1-dimensional spaces spanned by specific standard basis vectors in $V$, we would like $\phi(\wrap)$ to have distinct eigenvalues on the $v_k$. Furthermore, we would like to find a set of diagrammatic relations in the annular web category that enforces a choice of $\phi(\wrap)$ with distinct eigenvalues, or in other words, with a separable characteristic polynomial $\prod_{i=1}^N(X-\gamma_i)= \sum_{k=0}^N X^{N-k} (-1)^{k} e_k(\vec{\gamma})$. Here $e_k(\vec{\gamma})=e_k(\{\gamma_1,\dots,\gamma_N\})$ denotes the $k$-th elementary symmetric polynomial evaluated at the complex numbers $\gamma_1,\dots,\gamma_N$.

\begin{lemma} \label{lem:essentialideal} Suppose that $\phi$ is a functor from
$\AWeb[N]$ to complex vector spaces, which agrees on the subcategory $\Web[N]$
with the functor from Theorem~\ref{thm:CKM} and such that $\phi(\wrap)$ respects
the weight space decomposition of $V$. Then the coefficients of the
characteristic polynomial of $\phi(\wrap)$ are determined by the image of $\phi$
on endomorphisms of the empty object. More precisely:

 \begin{equation}
\label{eqn:essentialideal}
  \phi\left(
  \begin{tikzpicture}[anchorbase, scale=.3]
    %% Inner boundary circle
    \draw (0,0) circle (1);
    %% Outer boundary circle
    \draw (0,0) circle (3);
    %% Essential circle
    \draw [very thick, directed=.55] (0,0) circle (2);
    %% identity component
    % boundary markings
    \node at (0,-1) {$*$};
    \node at (-2.5,0) {\tiny$k$};
    \node at (0,-3) {$*$};
    \draw [dashed] (0,-1) to (0,-3);
  \end{tikzpicture}
  \right)
  = \sum_{|S|=k} (\prod_{i\in S} \gamma_i) \id_\C = e_k(\vec{\gamma})
\end{equation}
The annular web shown on the left-hand side will be called an \textit{essential
circle} in the annulus.
\end{lemma}



\begin{proof} The morphism can be written as the composition of a $k$-cup, a $k$-wrap and a $k$-cap. The sum $\sum_{|S|=k}$ comes from the cup and the factors from the action of the wrap on $v_S$.
\end{proof}

We now prescribe $\phi(\wrap)\colon V\to V$ to have the separable characteristic
polynomial $X^N-1$, and we may index the roots as $\gamma_k=\zeta^k=e^{k 2 \pi i
/N}$. This choice of relation is homogeneous with respect to a $\Z/N\/Z$-grading
by winding number, see Definition~\ref{def:grading}, and it has the effect of
setting $e_k(\vec{\gamma})=0$ for $1\leq k <N$ and
$e_N(\vec{\gamma})=(-1)^{N-1}$ in \eqref{eqn:essentialideal}. (More general evaluations are studied in \cite{LTV}.) We now extend the
definition of $\phi$ to the new generators $\wrap$ in the general case, that is,
allowing more than one strand.
\begin{defi}
  Let $V\otimes W$ be the image under $\phi$ of the domain of $\wrap$ and
  $W\otimes V$ its co-domain. Then we define $\phi(\wrap)$ to be the linear map
  determined by $v_{k}\otimes w\mapsto \zeta^{k} w\otimes v_k
  $
  for $v_{k}\in V$ and any $w\in W$. Furthermore we set $\phi(\wrapi)=\phi(\wrap)^{-1}$ and analogously for the duals.
\end{defi}

Let $\h$ denote the Cartan subalgebra of diagonal matrices in $\glnn{N}$,
$h_1,\dots,h_N$ its standard basis elements, and consider $U(\h)\subset
U(\glnn{N})$. We denote by $\reph$ the category of finite-dimensional
$U(\h)$-representations of integral weights, i.e. where the $h_i$ act semisimply
with integral eigenvalues. Note that the inclusion $\h\hookrightarrow \glnn{N}$
induces a restriction functor $\Rep(\glnn{N})\to\reph$ and that $\phi(\wrap)$
and $\phi(\wrapi)$ are morphisms in $\reph$. 


\begin{lemma}\label{lem:welldef} The functor $\phi\colon \AWeb[N] \to \reph$ is well-defined. 
\end{lemma}
\begin{proof} We will show that $\phi$ respects all relations defining the
morphism spaces of $\AWeb[N]$. All morphisms in $\AWeb[N]$ are linear
combinations of compositions of caps or cups between adjacent strands, vertices,
as well as the morphisms $\wrap$ and $\wrapi$, considered modulo local linear
relations and isotopy. As usual in skein theory, it suffices to impose the local
linear relations in a fixed disk, e.g. away from the seam $\alpha$. These
relations are respected by $\phi$ because we extended $\phi$ from
$\Web[N]\subset \AWeb[N]$. It thus remains to check that $\phi$ respects the
isotopy relations of sliding caps, cups and vertices along
wraps around the annulus. However, the images under $\phi$ of the wrap morphisms
were precisely chosen for these relations to hold (see Equation~\eqref{eq:capslide}). 
\end{proof}



\subsection{The tensor product on annular webs}
Let $\otimes\colon \AWeb[N]\times \AWeb[N]\to \AWeb[N]$ denote the bi-functor
given on objects by concatenation and on morphisms by superimposing a pair of
annular webs $(W_1,W_2)$ as follows:

\[
 \begin{tikzpicture}[anchorbase, scale=.30]
%% stuff inside -- inner
\draw[thick] (0,0) circle (2.5);
\fill[black,opacity=.2] (0,0) circle (2.5);
\draw[thick,fill=white] (0,0) circle (1.5);
%% stuff inside -- outer
\draw[thick] (0,0) circle (4.5);
\fill[black,opacity=.2] (4.5,0) arc (0:360:4.5) -- (3.5,0) arc (360:0:3.5);
\draw [thick] (0,0) circle (3.5);
%% Inner boundary circle
\draw (0,0) circle (1);
%% Outer boundary circle
\draw (0,0) circle (5);
%% left side
% \draw[dotted] (-2.29,2.29) to [out=225,in=90] (-3.25,0) to [out=270,in=135] (-2.29,-2.29);
% \draw[dotted] (-3.4,3.4) to [out=225,in=90] (-4.75,0) to [out=270,in=135] (-3.4,-3.4);
% \draw[dotted] (-1.93,1.93) to [out=225,in=90] (-2.75,0) to [out=270,in=135] (-1.93,-1.93);
% \draw[dotted] (-0.88,0.88) to [out=225,in=90] (-1.25,0) to [out=270,in=135] (-0.88,-0.88);
\draw [white,line width=.15cm] (-.6,.8) to (-1.8,2.4);
\draw [white,line width=.15cm] (-.8,.6) to (-2.4,1.8);
\draw [white,line width=.15cm] (-.95,.3) to (-2.85,0.9);
\draw [white,line width=.15cm] (-1,0) to (-3,0);
\draw [white,line width=.15cm] (-.95,-.3) to (-2.85,-.9);
\draw [white,line width=.15cm] (-.8,-.6) to (-2.4,-1.8);
\draw [white,line width=.15cm] (-.6,-.8) to (-1.8,-2.4);
\draw [very thick] (-.6,.8) to (-2.1,2.8);
\draw [very thick] (-.8,.6) to (-2.8,2.1);
\draw [very thick] (-.95,.3) to (-3.35,1.05);
\draw [very thick] (-1,0) to (-3.5,0);
\draw [very thick] (-.95,-.3) to (-3.35,-1.05);
\draw [very thick] (-.8,-.6) to (-2.8,-2.1);
\draw [very thick] (-.6,-.8) to (-2.1,-2.8);
\draw [very thick] (-2.7,3.6) to (-3,4);
\draw [very thick] (-3.6,2.7) to (-4,3);
\draw [very thick] (-4.25,1.35) to (-4.75,1.5);
\draw [very thick] (-4.5,0) -- (-5,0);
\draw [very thick] (-4.25,-1.35) to (-4.75,-1.5);
\draw [very thick] (-3.6,-2.7) to (-4,-3);
\draw [very thick] (-2.7,-3.6) to (-3,-4);
%%%%%right side
% \draw[dotted] (2.29,2.29) to [out=315,in=90] (3.25,0) to [out=270,in=45] (2.29,-2.29);
% \draw[dotted] (3.4,3.4) to [out=315,in=90] (4.75,0) to [out=270,in=45] (3.4,-3.4);
% \draw[dotted] (1.93,1.93) to [out=315,in=90] (2.75,0) to [out=270,in=45] (1.93,-1.93);
% \draw[dotted] (0.88,0.88) to [out=315,in=90] (1.25,0) to [out=270,in=45] (0.88,-0.88);
\draw [very thick] (1.5,2) to (1.98,2.64);
\draw [very thick] (2.79,3.72) to (3,4);
\draw [very thick] (2,1.5) to (2.64,1.98);
\draw [very thick] (3.72,2.79) to (4,3);
\draw [very thick] (2.375,.75) to (3.135,0.99);
\draw [very thick] (4.41,1.395) to (4.75,1.5);
\draw [very thick] (2.5,0) to (3.3,0);
\draw [very thick] (4.64,0) to (5,0);
\draw [very thick] (2.375,-.75) to (3.135,-0.99);
\draw [very thick] (4.41,-1.395) to (4.75,-1.5);
\draw [very thick] (2,-1.5) to (2.64,-1.98);
\draw [very thick] (3.72,-2.79) to (4,-3);
\draw [very thick] (1.5,-2) to (1.98,-2.64);
\draw [very thick] (2.79,-3.72) to (3,-4);
% inner
\draw [very thick] (.6,-.8) to (.9,-1.2);
\draw [very thick] (.8,-.6) to (1.2,-.9);
\draw [very thick] (.95,-.3) to (1.42,-.45);
\draw [very thick] (1,0) to (1.5,0);
\draw [very thick] (.95,.3) to (1.42,.45);
\draw [very thick] (.8,.6) to (1.2,.9);
\draw [very thick] (.6,.8) to (.9,1.2);
% boundary markings
\node at (0,-1) {$*$};
\node at (0,-5) {$*$};
\draw [dashed] (0,-1) to (0,-5);
% T_m
\node at (0,1.95) {\tiny$W_2$};
\node at (0,3.95) {\tiny$W_1$};
\end{tikzpicture}
\]
and resolving all crossings via \eqref{eq:crossing}. In words, we stack the
annular web $W_1$ on top of the annular web $W_2$, making sure that the ``legs'' of $W_1$
meet the inner (source) boundary of the annulus to the left of the legs of
$W_2$. Similarly, the ``arms'' of $W_1$ meet the outside (target) boundary of
the annulus to the left of all arms of $W_2$. This operation creates new
crossings, as the legs of $W_1$ might cross strings from $W_2$, and arms of
$W_2$ might cross strings from $W_1$. The operation is well-defined thanks to
Lemma \ref{lem:Reid}. In the above picture we have drawn $W_1$ stacked on top of
$W_2$, with strands leaving the former crossing over the latter. However, since
we are working at $q=1$, over- and under-crossings are equal and using isotopy
relations it is straightforward to check that $\otimes$ together with the
symmetric braiding from $\Web[N]$ yield a symmetric monoidal structure on
$\AWeb[N]$ such that both the inclusion $\Web[N]\to \AWeb[N]$ and $\phi \colon
\AWeb[N] \to \reph$ become symmetric monoidal functors. This also justifies why
we have described the action $\otimes$ to be ``superimposing'' annular webs. 

\begin{defi} \label{def:si} For $m\geq 0$ we denote by $\AWeb[N](m)$ the
endomorphism algebra in $\AWeb[N]$ of the object consisting of a sequence of $m$
points with label $1$ and outward orientation. We denote by $s_i$ for $1\leq i
\leq m$ the element of this endomorphism algebra that is given by the crossing
between the strands in positions $i$ and $i+1$, with positions understood modulo
$m$. We also write $u_i=\id_2-s_i$ for the corresponding \textit{dumbbell web}.

More generally, for two objects $\vec{k}$ and $\vec{l}$ we use the shorthand $\AWeb[N](\vec{k},\vec{l}):= \Hom_{\AWeb[N]}(\vec{k},\vec{l})$. We will also use these notation conventions for other categories. 
\end{defi}

The following lemma will allow us to freely express webs in terms of images of $1$-labeled tangles, which will be very useful in a number of proofs. 

\begin{lemma} \label{lem:resolveastangles}Every element of $\AWeb[N](m)$ for
$m\geq 0$ can be written as a $\C$-linear combination of $1$-labeled annular
tangles.
\end{lemma}
As we will see later, we may assume that the closed components are essential
circles (possibly carrying higher labels) and the non-closed components are
oriented $1$-labeled arcs from the inner to the outer boundary circle, that are
everywhere outward pointing. 
\begin{proof} We include this classical proof for completeness. It is well-known that every closed web can be written as a $\C$-linear combination of webs $W$ with only $1$-labeled edges, interacting at most in $2$-labeled dumbbells, see e.g. \cite[Proof of Lemma 4.1]{TVW}. Indeed, the argument can be inductively built from the following two operations:
\[
\begin{tikzpicture}[anchorbase]
  \draw [very thick] (0,0) -- (0,2);
  \node at (.3,1) {\tiny $k$};
\end{tikzpicture}
\quad = \frac{1}{k!}
\begin{tikzpicture}[anchorbase]
  \draw [very thick] (0,0) -- (0,.5);
  \draw [very thick] (0,.5) .. controls (-.5,.75) and (-.5,1.25) .. (0,1.5);
  \draw [very thick] (0,.5) .. controls (-.3,.75) and (-.3,1.25) .. (0,1.5);
  \node at (.15,1) {\small $\dots$};
  \draw [very thick] (0,.5) .. controls (.5,.75) and (.5,1.25) .. (0,1.5);
  \draw [very thick] (0,1.5) --(0,2);
\end{tikzpicture}
\quad \text{and}\quad
\begin{tikzpicture}[anchorbase]
  \draw [very thick] (-.5,0) .. controls (-.5,.25) .. (0,.5);
  \draw [very thick] (-.3,0) .. controls (-.3,.25) .. (0,.5);
  \node at (.1, .05) {\small $\dots$};
  \draw [very thick] (.5,0) .. controls (.5,.25) .. (0,.5);
  \draw [very thick] (0,.5) -- (0,1);
  \draw [very thick] (0,1) .. controls (-.5,1.25) .. (-.5,1.5);
  \node at (-.1,1.45) {\small $\dots$};
  \draw [very thick] (0,1) .. controls (.3,1.25) .. (.3,1.5);
  \draw [very thick] (0,1) .. controls (.5,1.25) .. (.5,1.5);
  \draw[decorate,decoration={brace}] (.5,-.1) -- node[midway,font=\small,yshift=-.25cm] {$k$} (-.5,-.1);
  \draw[decorate,decoration={brace}] (-.5,1.6) -- node[midway,font=\small,yshift=.25cm] {$k$} (.5,1.6);
\end{tikzpicture}
\quad = \;(-1)^k\;\;
\begin{tikzpicture}[anchorbase]
  \draw [very thick] (-.5,0) .. controls (-.5,.75) and (.5,.75) .. (.5,1.5);
  \draw [very thick] (-.3,0) .. controls (-.3,.25) .. (.1,.5);
  \node at (.1, .05) {\small $\dots$};
  \draw [very thick] (.5,0) .. controls (.5,.25) .. (.1,.5);
  \draw [very thick] (.1,.5) .. controls (.1,.75) and (-.1,.75) .. (-.1,1);
  \draw [very thick] (-.1,1) .. controls (-.5,1.25) .. (-.5,1.5);
  \node at (-.1,1.45) {\small $\dots$};
  \draw [very thick] (-.1,1) .. controls (.3,1.25) .. (.3,1.5);
  \draw[decorate,decoration={brace}] (.5,-.1) -- node[midway,font=\small,yshift=-.25cm] {$k-1$} (-.3,-.1);
  \draw[decorate,decoration={brace}] (-.5,1.6) -- node[midway,font=\small,yshift=.25cm] {$k-1$} (.3,1.6);
\end{tikzpicture}
\quad + \; (-1)^{k+1}\;\;
\begin{tikzpicture}[anchorbase]
  \draw [very thick] (-.5,0) .. controls (-.5,.25) and (-.2,.5) .. (-.2,.85);
  \draw [very thick] (-.2,.85) .. controls (-.2,.9) and (-.1,.9) .. (-.1,1);
  \draw [very thick] (-.3,0) .. controls (-.3,.25) .. (.1,.5);
  \node at (.1, .05) {\small $\dots$};
  \draw [very thick] (.5,0) .. controls (.5,.25) .. (.1,.5);
  \draw [very thick] (.1,.5) .. controls (.1,.6) and (.2,.6) .. (.2,.65);
  \draw [very thick] (.2,.65) .. controls (.1,.75) and (-.1,.75) .. (-.2,.85);
  \draw [very thick] (.2,.65) .. controls (.2,1) and (.5,1.25) .. (.5,1.5);
  \draw [very thick] (-.1,1) .. controls (-.5,1.25) .. (-.5,1.5);
  \node at (-.1,1.45) {\small $\dots$};
  \draw [very thick] (-.1,1) .. controls (.3,1.25) .. (.3,1.5);
  \draw[decorate,decoration={brace}] (.5,-.1) -- node[midway,font=\small,yshift=-.25cm] {$k-1$} (-.3,-.1);
  \draw[decorate,decoration={brace}] (-.5,1.6) -- node[midway,font=\small,yshift=.25cm] {$k-1$} (.3,1.6);
\end{tikzpicture}
\] 
The remaining $2$-labeled dumbbells can now be expanded in terms of crossings and their oriented resolutions, resulting in a linear combination of $1$-labeled annular tangles.
\end{proof}
An analogous result is true at generic $q$ over $\C(q)$.

\begin{lemma}\label{lem:endempty} The endomorphism algebra of the empty object
in $\AWeb[N]$ is isomorphic to $\C[c_1,\dots, c_{N-1}, c_N^{\pm 1}]$, where
$c_i$ denotes the counter-clockwise oriented $i$-labeled essential circle in the
annulus and $c_N^{-1}$ the $N$-labeled clockwise oriented essential circle.
\end{lemma}
\begin{proof}
The proof proceeds in two steps. First we show that any closed web in the
annulus can be written as a $\C$-linear combination of collections of essential
circles from the set $\{c_1,c_2, \dots, c_N, c_N^{-1}\}$. The essential circles
commute, which is easily seen by applying Reidemeister II moves, and $c_N$ and
$c_N^{-1}$ are indeed mutually inverse by \eqref{eq:webrel2}. Second, we check that the
counter-clockwise oriented circles are algebraically independent. 

  By Lemma~\ref{lem:resolveastangles} any closed web may be expressed as a linear
    combination of $1$-labeled annular links. Recall that we are working at $q=1$,
    so crossings signs are irrelevant and the components of the link can be pulled
    apart into a disjoint union. Moreover, each of the components can be assumed
    to wrap exclusively clockwise or counter-clockwise around the annulus. (By
    this we mean that upon choosing an orientation of the seam $\alpha$, all
    intersection points of $\alpha$ with the link component have the same sign.)
    The annular evaluation algorithm from \cite[Lemma 5.2]{QR2} implies that any
    such web can be resolved in terms of concentric essential circles of various
    labels. Hence, the same is true for the original web. To finish the first
    part, note that the clockwise-oriented circle of label $i$ is equal to
    $c_{N-i}c_N^{-1}$. This follows from the last relation in \eqref{eq:webrel2}. 




In order to prove algebraic independence, we use an extension of the functor
$\phi$ from Lemma~\ref{lem:welldef}. Let $R=\C[\X^{\pm 1} ]=\C[X_1^{\pm
1},\dots, X_N^{\pm 1}]$ be a Laurent polynomial ring in $N$ variables and
consider the category $R\reph$ obtained by tensoring all morphism spaces in
$\reph$ by $R$. Then $\phi_R\colon \AWeb[N] \to R\reph$ can be defined just as
$\phi$ was above, except that the $X_i$ now play the role of the eigenvalues of
the $1$-labeled counter-clockwise wrap: $\phi_R(D)(v_k):=X_k v_k$. 

For a closed annular web $W$, it follows that $\phi(W)\in R$. In particular,
counter-clockwise oriented circles $c_i$ evaluate to elementary symmetric
polynomials $e_i(\X)$ in the variables $X_k$ and their clockwise counterparts
evaluate to $e_i(\X^{-1})$. As a consequence of the first part of this proof,
the evaluation of closed annular webs $\phi_R(W)$ takes values in the symmetric
part $R^{\mathfrak{S}_N} \cong \C[e_1(\X),\dots, e_{N-1}(\X), e_{N}(\X)^{\pm
1}]$ of $R$. Now the algebraic independence of the $c_i$ follows from the
algebraic independence of their images $e_i(\X)$ under $\phi_R$.
\end{proof}

An analogous version of this result holds for $\AWebq[N]$ over $\C(q)$. The second part in its $q=1$ version is already sufficient to establish algebraic independence of the counter-clockwise essential circles in $\AWebq[N]$.

Just as in the non-annular case, we denote by $\AWebp[N]$ the full subcategory
of $\AWeb[N]$ with objects given by all upward, or \emph{outward}, pointing
boundary sequences. The following is the analog of Lemma~\ref{lem:upward}.

\begin{lemma}\label{lem:upwardannulus} The morphism spaces of $\AWebp[N]$ are
spanned by webs with all edges outward oriented (i.e. admitting an oriented parametrisation with
derivative having a positive radial component everywhere), but potentially superimposed with
essential circles.
\end{lemma}

Note that Lemma \ref{lem:endempty} allows us to restrict to counter-clockwise
essential circles, except for the $N$-labeled ones.

\begin{proof}
 We will prove the claim for webs $W$ whose source and target objects are only $1$-labeled. The general claim follows by a usual merging argument.

Lemma~\ref{lem:resolveastangles} allows us to write $W$ as a linear combination of annular tangles. The closed components of these tangles evaluate to essential circles thanks to Lemma~\ref{lem:endempty}, while the non-closed components can be isotoped to be outward oriented arcs. (This shortcut is possible since we work at $q=1$.) The superposition of such outward arcs is itself a linear combination of outward webs, as desired. 
\end{proof}




\subsection{Equivalences between blocks}\label{sec:lambdaequiv}


The categories $\Web[N]$, $\AWeb[N]$ and all of their specializations and
quotients considered in the following decompose into blocks (i.e. full
subcategories) indexed by $m\in \Z$, which consist of those objects, whose
signed sum of labels equals $m$. Here we count upward oriented boundary points
positively, and downward pointing ones negatively. We indicate such blocks by
the subscript $m$, e.g. $\Web[N]_{m}$.

We will use the notation $\lambda$ for the endofunctors of these categories that
act on objects by tensoring with a single $N$-labeled upward boundary point on
the right and on morphisms by superimposing with an upward-oriented $N$-labeled
edge parallel to the seam and with boundary points to the right of all other
boundary points. We denote by $\lambda^*$ the analogous operation with downward
orientations. We give an example for $\AWeb[N]$:
\[
\begin{tikzpicture}[anchorbase, scale=.4]
%% stuff inside 
\draw[thick] (0,0) circle (2.5);
\fill[black,opacity=.2] (0,0) circle (2.5);
\draw[thick,fill=white] (0,0) circle (1.5);
\draw[dotted] (-1.93,1.93) to [out=45,in=180] (0,2.75) to [out=0,in=135] (1.93,1.93);
\draw[dotted] (-0.88,0.88) to [out=45,in=180] (0,1.25) to [out=0,in=135] (0.88,0.88);
%% Inner boundary circle
\draw (0,0) circle (1);
%% Outer boundary circle
\draw (0,0) circle (3);
%% web edges
\draw [very thick] (.8,.6) to (1.2,.9);
\draw [very thick] (-.8,.6) to (-1.2,.9);
\draw [very thick] (2,1.5) to (2.4,1.8);
\draw [very thick] (-2,1.5) to (-2.4,1.8);
% boundary markings
\node at (0,-1) {$*$};
\node at (0,-3) {$*$};
\draw [dashed] (0,-1) to (0,-3);
\end{tikzpicture}
\quad
\xrightarrow{\lambda}
\quad 
 \begin{tikzpicture}[anchorbase, scale=.4]
%% stuff inside 
\draw[thick] (0,0) circle (2.5);
\fill[black,opacity=.2] (0,0) circle (2.5);
\draw[thick,fill=white] (0,0) circle (1.5);
\draw[dotted] (-1.93,1.93) to [out=45,in=180] (0,2.75) to [out=0,in=135] (1.93,1.93);
\draw[dotted] (-0.88,0.88) to [out=45,in=180] (0,1.25) to [out=0,in=135] (0.88,0.88);
%% web edges
\draw [very thick] (.8,.6) to (1.2,.9);
\draw [very thick] (-.8,.6) to (-1.2,.9);
\draw [very thick] (2,1.5) to (2.4,1.8);
\draw [very thick] (-2,1.5) to (-2.4,1.8);
\draw [white, line width=.12cm] (1,0) to (3,0);
\draw [very thick,->] (1,0) to (3,0);
\node at (3.3,0) {\tiny $N$};
% boundary markings
\node at (0,-1) {$*$};
\node at (0,-3) {$*$};
\draw [dashed] (0,-1) to (0,-3);
%% Inner boundary circle
\draw (0,0) circle (1);
%% Outer boundary circle
\draw (0,0) circle (3);
\end{tikzpicture}
\quad
\xrightarrow{\lambda^*}
\quad 
 \begin{tikzpicture}[anchorbase, scale=.4]
%% stuff inside 
\draw[thick] (0,0) circle (2.5);
\fill[black,opacity=.2] (0,0) circle (2.5);
\draw[thick,fill=white] (0,0) circle (1.5);
\draw[dotted] (-1.93,1.93) to [out=45,in=180] (0,2.75) to [out=0,in=135] (1.93,1.93);
\draw[dotted] (-0.88,0.88) to [out=45,in=180] (0,1.25) to [out=0,in=135] (0.88,0.88);
%% web edges
\draw [very thick] (.8,.6) to (1.2,.9);
\draw [very thick] (-.8,.6) to (-1.2,.9);
\draw [very thick] (2,1.5) to (2.4,1.8);
\draw [very thick] (-2,1.5) to (-2.4,1.8);
\draw [white, line width=.12cm] (1,0) to (3,0);
\draw [very thick,->] (1,0) to (3,0); 
\node at (3.3,0) {\tiny $N$};
\draw [white, line width=.12cm] (.8,-.6) to (2.4,-1.8);
\draw [very thick,<-] (.8,-.6) to (2.4,-1.8);
\node at (2.64,-1.98) {\tiny $N$};
% boundary markings
\node at (0,-1) {$*$};
\node at (0,-3) {$*$};
\draw [dashed] (0,-1) to (0,-3);
%% Inner boundary circle
\draw (0,0) circle (1);
%% Outer boundary circle
\draw (0,0) circle (3);
\end{tikzpicture}
\] 
Here we display the $N$-labeled strand as crossing over the remaining web for better visibility, even though this has no significance at $q=1$.

\begin{lemma} The endofunctors $\lambda$ and $\lambda^*$ establish an
equivalence between the blocks $\AWeb[N]_m$ and $\AWeb[N]_{m+N}$.
\end{lemma}
\begin{proof}
Using isotopy relations, it is easy to see that the following types of webs provide natural isomorphisms between the identity functor on $\AWeb[N]_m$ and the endofunctor given by the composition $\lambda^*\circ \lambda$:
\[ \begin{tikzpicture}[anchorbase, scale=.4]
%% stuff inside 
\draw[dotted] (-1.93,1.93) to [out=45,in=180] (0,2.75) to [out=0,in=135] (1.93,1.93);
\draw[dotted] (-0.88,0.88) to [out=45,in=180] (0,1.25) to [out=0,in=135] (0.88,0.88);
%% web edges
\draw [very thick] (.8,.6) to (2.4,1.8);
\draw [very thick] (-.8,.6) to (-2.4,1.8);
\draw [white, line width=.12cm] (2.4,-1.8) to (1.2,-.9) to  [out=135, in=180] (1.5,0)to (3,0);
\draw [very thick,->] (2.4,-1.8) to (1.2,-.9) to  [out=135, in=180] (1.5,0)to (3,0);
% boundary markings
\node at (3.3,0) {\tiny $N$};
\node at (0,-1) {$*$};
\node at (0,-3) {$*$};
\draw [dashed] (0,-1) to (0,-3);
%% Inner boundary circle
\draw (0,0) circle (1);
%% Outer boundary circle
\draw (0,0) circle (3);
\end{tikzpicture}
\] 
Analogously, there are natural transformations between $\lambda \circ \lambda^*$ and the identity functor on $\AWeb[N]_{m}$. 
\end{proof}

\subsection{The quotients by essential circles}
It is a key observation that the functor $\phi \colon \AWeb[N] \to \reph$ is not faithful.
\begin{prop}\label{prop:nff} We have the following identities in $\reph$:
 \begin{equation}
\label{eqn:essentialidealtwo}
  \phi\left(
  \begin{tikzpicture}[anchorbase, scale=.3]
    %% Inner boundary circle
    \draw (0,0) circle (1);
    %% Outer boundary circle
    \draw (0,0) circle (3);
    %% Essential circle
    \draw [very thick,directed=.55] (0,0) circle (2);
    %% identity component
  \draw [very thick] (-.9,.4) to (-2.8,1.2);
    \node at (0,1.35) {$\cdots$};
    \node at (0,2.5) {$\cdots$};
    \draw [very thick] (.9,.4) to (2.8,1.2);
    % boundary markings
    \node at (0,-1) {$*$};
    \node at (0,-3) {$*$};
    \draw [dashed] (0,-1) to (0,-3);
    \node at (-2.5,0) {\tiny $k$};
  \end{tikzpicture}
  \right)
  =
   \phi\left(
  \begin{tikzpicture}[anchorbase, scale=.3]
    %% Inner boundary circle
    \draw (0,0) circle (1);
    %% Outer boundary circle
    \draw (0,0) circle (3);
    %% Essential circle
    \draw [very thick,rdirected=.55] (0,0) circle (2);
    %% identity component
  \draw [very thick] (-.9,.4) to (-2.8,1.2);
    \node at (0,1.35) {$\cdots$};
    \node at (0,2.5) {$\cdots$};
    \draw [very thick] (.9,.4) to (2.8,1.2);
    % boundary markings
    \node at (0,-1) {$*$};
    \node at (0,-3) {$*$};
    \draw [dashed] (0,-1) to (0,-3);
    \node at (-2.5,0) {\tiny $k$};
  \end{tikzpicture}
  \right)
  =\begin{cases} (-1)^{N-1} \id & \text{ if } k=N\\
  0 & \text{ otherwise }
  \end{cases}.
\end{equation}
\end{prop}
\begin{proof}
  % This follows from \eqref{eqn:essentialideal}.
 First consider the case where there are no strands mapping between the inner and outer circles. In this case, the result follows by specializing~\eqref{eqn:essentialideal} at $\gamma_l=e^{k 2 \pi i
   /N}$.

 In the case where there are strands crossing the essential $k$-labeled circle,
 the result follows similarly: a crossing is sent under $\phi$ to the
 transposition $u\otimes v \rightarrow v\otimes u$. Then the computation in the
 proof of Lemma~\ref{lem:essentialideal} can be performed even in the presence
 of extra strands, yielding the same scalar.
\end{proof}


\begin{defi} We let $\essbAWeb[N]$ denote the quotient of $\AWeb[N]$ by
the tensor ideal generated by the $k$-labeled essential circles $c_k$ for $k<N$. We
define $\essAWeb[N]$ to be quotient of $\AWeb[N]$ by the tensor ideal generated
by the $c_k$ for $k<N$ as well as $c_N+(-1)^N$. 
\end{defi} 

Note that the monoidal structure $\otimes$ and the functor $\phi$ descend to the
quotients $\essbAWeb[N]$ and $\essAWeb[N]$. One of the key results of this paper
will be to prove that this category $\essAWeb[N]$ is equivalent to the full
subcategory of $\reph$ generated by the images of the objects of
$\Rep(\glnn{N})$.

Upon reinterpreting linear categories as locally unital algebras, $\essbAWeb[N]$
can be thought of as a central extension of $\essAWeb[N]$. For every object
there exists an automorphism, always denoted $c_N$, given by the $N$-labeled
counter-clockwise oriented essential circle superimposed onto the identity
morphism of the object, which ``commutes'' with all morphisms. Then
$\essAWeb[N]$ is obtained from $\essbAWeb[N]$ by setting $c_N=(-1)^{N-1}$.

\begin{defi}\label{def:grading} Let $W$ be a web in $\AWeb[N]$. Then the
\emph{flow winding number} $w_f(W)$ of $W$ is given by the algebraic
intersection number of the web with the segment $\alpha$ (assuming no trivalent
vertex occurs on it), where $k$-labeled edges crossing $\alpha$ count as $\pm
k$.
\end{defi}
It is clear that all web relations in $\AWeb[N]$ and $\essbAWeb[N]$ preserve the
flow winding number. This $\Z$-grading on the morphism spaces of these
categories descends to a $\Z/N\Z$ grading on $\essAWeb[N]$.

The following corollaries are implied by Lemma~\ref{lem:endempty} and
Lemma~\ref{lem:upwardannulus}.


\begin{coro}\label{cor:end} The endomorphism algebra of the empty object in
$\essAWeb[N]$ is $\C$ and in $\essbAWeb[N]$ it is isomorphic to the Laurent
polynomial ring over $\C$ generated by an essential $N$-labeled circle. 
\end{coro}

\begin{coro}\label{cor:out} The morphism spaces of $\essAWebp[N]$ are
spanned by outward pointing webs. The same is true in $\essbAWebp[N]$ up to
superposition with an integer power of the $N$-labeled essential circle $c_N$.
\end{coro}



\begin{lemma} In $\essbAWeb[N](1)$ we have $\wrap^N=(-1)^{N-1}c_N$, where we
abuse notation by writing $c_N$ for the superposition of the identity on $1$
with the counterclockwise $N$-labeled essential circle. In $\essAWeb[N](1)$ this
specializes to $\wrap^N = \id_1$.
\label{lem:DNone}
\end{lemma}
\begin{proof} We compute:
 \begin{align*} 
 (-1)^{N-1} c_N
  \; &= (-1)^{N-1}\; \begin{tikzpicture}[anchorbase, scale=.3]
    %% Inner boundary circle
    \draw (0,0) circle (1);
    %% Outer boundary circle
    \draw (0,0) circle (3);
    \draw [very thick,->] (0,1) to (0,3);
    %% Essential circle
    \draw [very thick,directed=.55] (0,0) circle (2);
    % boundary markings
    \node at (0,-1) {$*$};
    \node at (0,-3) {$*$};
        \node at (0.5,-2.4) {\tiny $N$};
    \draw [dashed] (0,-1) to (0,-3);
  \end{tikzpicture}
  \; = \; \begin{tikzpicture}[anchorbase, scale=.3]
    %% Inner boundary circle
    \draw (0,0) circle (1);
    %% Outer boundary circle
    \draw (0,0) circle (3);
    \draw [very thick,directed=.65] (0,1) to [out=90,in=270] (-.8,1.8);
     \draw [very thick,->] (.8,1.8) to [out=90,in=270] (0,2.8) to (0,3);
    %% Essential circle
    \draw [very thick,directed=.55,directed=.25] (0,0) circle (2);
    % boundary markings
    \node at (0,-1) {$*$};
    \node at (0,-3) {$*$};
            \node at (0.5,-2.4) {\tiny $N$};
    \draw [dashed] (0,-1) to (0,-3);
  \end{tikzpicture}
    \; = \; \begin{tikzpicture}[anchorbase, scale=.3]
    %% Inner boundary circle
    \draw (0,0) circle (1);
    %% Outer boundary circle
    \draw (0,0) circle (3);
    \draw [very thick,directed=.65] (0,1) to [out=90,in=270] (.8,1.8);
     \draw [very thick,->] (-.8,1.8) to [out=90,in=270] (0,2.8) to (0,3);
    %% Essential circle
    \draw [very thick,directed=.55,directed=.25] (0,0) circle (2);
    % boundary markings
    \node at (0,-1) {$*$};
    \node at (0,-3) {$*$};
            \node at (.8,-2.4) {\tiny $N$-$1$};
    \draw [dashed] (0,-1) to (0,-3);
  \end{tikzpicture} \circ \wrap
  \\
  &= (-1)^{N-1} \; \begin{tikzpicture}[anchorbase, scale=.3]
    %% Inner boundary circle
    \draw (0,0) circle (1);
    %% Outer boundary circle
    \draw (0,0) circle (3);
    \draw [very thick,->] (0,1) to (0,3);
    %% Essential circle
    \draw [very thick,directed=.55] (0,0) circle (2);
    % boundary markings
    \node at (0,-1) {$*$};
    \node at (0,-3) {$*$};
        \node at (0.8,-2.4) {\tiny $N$-$1$};
    \draw [dashed] (0,-1) to (0,-3);
  \end{tikzpicture} \circ \wrap
  \;+\; 
  \begin{tikzpicture}[anchorbase, scale=.3]
    %% Inner boundary circle
    \draw (0,0) circle (1);
    %% Outer boundary circle
    \draw (0,0) circle (3);
    \draw [very thick,directed=.65] (0,1) to [out=90,in=270] (-.8,1.8);
     \draw [very thick,->] (.8,1.8) to [out=90,in=270] (0,2.8) to (0,3);
    %% Essential circle
    \draw [very thick,directed=.55,directed=.25] (0,0) circle (2);
    %% identity component
  %\draw [very thick] (-.9,.4) to (-2.8,1.2);
  %  \draw [very thick] (.9,.4) to (2.8,1.2);
    % boundary markings
    \node at (0,-1) {$*$};
    \node at (0,-3) {$*$};
            \node at (0.8,-2.4) {\tiny $N$-$1$};
    \draw [dashed] (0,-1) to (0,-3);
  \end{tikzpicture}
  \circ \wrap
   \; = \; \begin{tikzpicture}[anchorbase, scale=.3]
    %% Inner boundary circle
    \draw (0,0) circle (1);
    %% Outer boundary circle
    \draw (0,0) circle (3);
    \draw [very thick,directed=.65] (0,1) to [out=90,in=270] (.8,1.8);
     \draw [very thick,->] (-.8,1.8) to [out=90,in=270] (0,2.8) to (0,3);
    %% Essential circle
    \draw [very thick,directed=.55,directed=.25] (0,0) circle (2);
    % boundary markings
    \node at (0,-1) {$*$};
    \node at (0,-3) {$*$};
            \node at (.8,-2.4) {\tiny $N$-$2$};
    \draw [dashed] (0,-1) to (0,-3);
  \end{tikzpicture} \circ \wrap^2
  \\
  &= \cdots = - \; \begin{tikzpicture}[anchorbase, scale=.3]
    %% Inner boundary circle
    \draw (0,0) circle (1);
    %% Outer boundary circle
    \draw (0,0) circle (3);
    \draw [very thick,->] (0,1) to (0,3);
    %% Essential circle
    \draw [very thick,directed=.55] (0,0) circle (2);
    % boundary markings
    \node at (0,-1) {$*$};
    \node at (0,-3) {$*$};
        \node at (0.5,-2.4) {\tiny $1$};
    \draw [dashed] (0,-1) to (0,-3);
  \end{tikzpicture} \circ \wrap^{N-1} 
  \;+\; \wrap^N = \wrap^N
 \end{align*}
The second equality arises by resolving the crossing while the third is an
isotopy. The equalities then alternate between such which hold by expanding a
crossing and such that use isotopies and \eqref{eqn:essentialidealtwo}.
\end{proof}
\begin{rema}An analogous argument shows $\sum_{i=0}^N D^{N-i} (-1)^i c_i =0$ in
$\AWeb[N]$, c.f. \cite[Section 8.2]{CK_ann}.
\end{rema}

Consider the algebra $\C[\wrap^{\pm 1}]/\la \wrap^N-1\ra$. Lemma~\ref{lem:DNone}
implies that this surjects onto the subalgebra of $\essAWeb[N](1)$ generated by
the wrapping morphisms $D$ and $D^{-1}$, and the flow winding grading implies that the surjection is an
isomorphism. The Chinese remainder theorem implies $\C[\wrap^{\pm 1}]/\la
\wrap^N-1\ra\cong \bigoplus_{j=1}^N \C[\wrap]/\la \wrap-e^{j 2 \pi i /N}\ra$.

\begin{defi} \label{def:Pk} For $1\leq k\leq N$ we denote by $P_k\in
\C[\wrap]$ a chosen representative for the idempotent that projects onto the
direct summand $\C[\wrap]/\la \wrap-e^{k 2 \pi i /N}\ra$ of $\C[\wrap^{\pm
1}]/\la \wrap^N-1\ra$.
\end{defi}
By abuse of notation we also write $P_k$ for the corresponding orthogonal idempotents in
$\essAWeb[N](1)$. It is a straightforward but crucial observation that
$\phi(P_k(\wrap))$ is the projection $V\twoheadrightarrow \C\la v_k\ra
\hookrightarrow V$. 


\begin{theo} \label{thm:fullness}
  The functor $\phi\colon \AWeb[N]\to \reph$ is full. 
\end{theo}


\begin{proof} We show that the induced functor on the quotient $\essAWeb[N]$ is
full. For this, let $\vec{k}$ and $\vec{l}$ be objects in $\essAWeb[N]$ and
$\cev{k}$ be the dual of $\vec{k}$, which is obtained from $\vec{k}$ by inverting orientations and the order of
the sequence. We consider the duality isomorphism $f\colon
\essAWeb[N](\vec{k},\vec{l}) \to \essAWeb[N](\emptyset,\vec{l} \otimes \cev{k})$
and its inverse, which can be explicitly described as the following operations on diagrams.
\[  f\colon 
  \begin{tikzpicture}[anchorbase, scale=.3]
    %% stuff inside -- inner
    \draw[thick] (0,0) circle (3.5);
    \fill[black,opacity=.2] (0,0) circle (3.5);
    \draw[thick,fill=white] (0,0) circle (2.5);
    %% Inner boundary circle
    \draw (0,0) circle (1);
    %% Outer boundary circle
    \draw (0,0) circle (4);
    %% left side
    \draw[dotted] (-2.65,2.65) to [out=225,in=90] (-3.75,0) to [out=270,in=135] (-2.65,-2.65);
    \draw[dotted] (-1.59,1.59) to [out=225,in=90] (-2.25,0) to [out=270,in=135] (-1.59,-1.59);
    \draw [very thick] (-2.4,3.2) to (-2.1,2.8);
    \draw [very thick] (-2.4,-3.2) to (-2.1,-2.8);
    %%%%%right side
    \draw[very thick] (-1.5,2) to (-0.6,0.8);
    \draw[very thick] (-1.5,-2) to (-0.6,-0.8);
    % boundary markings
    \node at (0,-1) {$*$};
    \node at (0,-4) {$*$};
    \draw [dashed] (0,-1) to (0,-4);
    %\node at (0,2.95) {$W$};
    \end{tikzpicture} 
    \mapsto 
\begin{tikzpicture}[anchorbase, scale=.3]
%% stuff inside -- inner
\draw[thick] (0,0) circle (3.5);
\fill[black,opacity=.2] (0,0) circle (3.5);
\draw[thick,fill=white] (0,0) circle (2.5);
%% Inner boundary circle
\draw (0,0) circle (1);
%% Outer boundary circle
\draw (0,0) circle (4);
%% left side
\draw[dotted] (-2.65,2.65) to [out=225,in=90] (-3.75,0) to [out=270,in=135] (-2.65,-2.65);
\draw[dotted] (-1.59,1.59) to [out=225,in=90] (-2.25,0) to [out=270,in=135] (-1.59,-1.59);
\draw [very thick] (-2.4,3.2) to (-2.1,2.8);
\draw [very thick] (-2.4,-3.2) to (-2.1,-2.8);
%%%%%right side
\draw [white,line width=.15cm] (0,2.4) to (0,3.6);
\draw[very thick] (-1.5,2) to [out=300,in=270] (0,2.5) to (0,4);
\draw [white,line width=.15cm] (1.97,1.47) to (3.16,2.37);
\draw[very thick] (-1.5,-2) to [out=60,in=270] (-1.75,0) to [out=90,in=180] (0,1.75) to [out=0,in=210] (2,1.5) to (3.2,2.4);
% boundary markings
\node at (0,-1) {$*$};
\node at (0,-4) {$*$};
\draw [dashed] (0,-1) to (0,-4);
%\node at (0,2.95) {$W$};
\end{tikzpicture}  
\quad,\quad 
f^{-1}\colon 
\begin{tikzpicture}[anchorbase, scale=.3]
  %% stuff inside -- inner
  \draw[thick] (0,0) circle (2.5);
  \fill[black,opacity=.2] (0,0) circle (2.5);
  \draw[thick,fill=white] (0,0) circle (1.5);
  %% Inner boundary circle
  \draw (0,0) circle (1);
  %% Outer boundary circle
  \draw (0,0) circle (4);
  %% left side
  \draw[dotted] (-2.29,2.29) to [out=225,in=90] (-3.25,0) to [out=270,in=135] (-2.29,-2.29);
  \draw [very thick] (-2.4,3.2) to (-1.5,2);
  \draw [very thick] (-2.4,-3.2) to (-1.5,-2);
  %%%%%right side
  \draw [very thick] (1.5,2) to (2.4,3.2);
  \draw [very thick] (0,4) to (0,2.5);
  \draw[dotted] (1.95,2.6) to [out=135,in=0] (0,3.25) ;
  % boundary markings
  \node at (0,-1) {$*$};
  \node at (0,-4) {$*$};
  \draw [dashed] (0,-1) to (0,-4);
  \end{tikzpicture}
\mapsto
 \begin{tikzpicture}[anchorbase, scale=.3]
%% stuff inside -- inner
\draw[thick] (0,0) circle (2.5);
\fill[black,opacity=.2] (0,0) circle (2.5);
\draw[thick,fill=white] (0,0) circle (1.5);
%% Inner boundary circle
\draw (0,0) circle (1);
%% Outer boundary circle
\draw (0,0) circle (4);
%% left side
\draw[dotted] (-2.29,2.29) to [out=225,in=90] (-3.25,0) to [out=270,in=135] (-2.29,-2.29);
\draw [very thick] (-2.4,3.2) to (-1.5,2);
\draw [very thick] (-2.4,-3.2) to (-1.5,-2);
%%%%%right side
\draw [white,line width=.15cm] (1,0) to (2.6,0);
\draw [white,line width=.15cm] (.6,-.8) to (1.8,-2.4);
\draw [very thick] (1,0) to (2.5,0) to [out=0,in=30] (1.5,2);
\draw [very thick] (.6,-.8) to (1.8,-2.4) to[out=300,in=225]  (2.46,-2.46) to [out=45,in=285] (3.38,.9) to[out=105,in=345] (.9,3.38) to [out=165,in=90] (0,3) to (0,2.5);
\draw[dotted] (1.94,1.94) to [out=135,in=0] (0,2.75) ;
\draw[dotted] (1.94,-1.94) to [out=45,in=270] (2.75,0) ;
\draw[dotted]  (1.25,0) to [out=270,in=45] (0.88,-0.88);
% boundary markings
\node at (0,-1) {$*$};
\node at (0,-4) {$*$};
\draw [dashed] (0,-1) to (0,-4);
\end{tikzpicture}
\] Since these operations are given by tensoring with an identity morphism and
then pre-composing with cups, or post-composing with caps, there are
corresponding isomorphisms $\phi(f^{\pm 1})$ between the relevant morphism
spaces in the target category $\reph$. 

To prove the theorem, it now suffices to
check that $\phi$ restricts to a surjective map from the morphism space
$\essAWeb[N](\emptyset,\vec{l} \otimes \cev{k})$ to
$\reph(\C,\phi(\vec{l}\otimes \cev{k}))$, the latter of which is
isomorphic to the zero weight space in $\phi(\vec{l}\otimes \cev{k})$ via the
map that evaluates elements of the morphism space at $1\in \C$. 

    Next we claim that it is enough to prove surjectivity in the case where
    $\vec{l}\otimes \cev{k}$ consists entirely of entries $1$, the first $n$ of
    which point outward and the last $n$ inward. Indeed, if all weight zero
    vectors in this tensor product are in the image (upon evaluation at $1\in
    \C$), then composing with the images of merge webs and the symmetric
    braiding, one may find any weight zero vector of a tensor product of
    fundamentals and their duals in the respective image.
    
Actually, it is sufficient to find for every $1\leq k\leq N$ a morphism $L_k\in
\essAWeb[N](\emptyset,(1,1^*))$, such $v_k\otimes v_k^*=\phi(L_k)(1)$, as then
we can take diagrammatic tensor products of such morphisms $L_{k_i}$, composed
with permutations, to find any standard basis vector of weight zero in the
image. To show that $v_k\otimes v_k^*$ is in the image, we invert the bending
process via the isomorphisms $f^{\pm 1}$ and $\phi(f^{\pm 1})$, and the problem
becomes equivalent to finding the projection $V \twoheadrightarrow \C\la
v_k\ra\hookrightarrow V$ in the image of $\phi$. But this we have already seen;
it is given by $\phi(P_k)$.
\end{proof}


Later, we will prove that $\phi$ induces a functor from $\essAWeb[N]$ to $\reph$
that is not only full, but also faithful (see Theorem \ref{thm:faithfulness}).
Nevertheless, we will continue to work in the more general framework of
$\essbAWeb[N]$ whenever possible.



%---------------------------------------------------------------------------------------
\subsection{Extremal weight projectors}
\label{sec:EWP}
%---------------------------------------------------------------------------------------
In \cite{QW}, we defined the concept of extremal weight projectors in the
context of (affine) $\slnn{2}$ skein theory. This involved finding a suitable
quotient of the affine Temperley-Lieb category, in which we identified a family
of idempotents akin to Jones-Wenzl projectors and corresponding, on the
representation-theoretic side, to projections onto the direct sum of the top and
bottom weight spaces in the tensor powers of the vector representation of
$U(\slnn{2})$. The same question naturally extends beyond the $\slnn{2}$ case,
and the definition can be adapted and generalized to the $\glnn{N}$ case as
follows.

\begin{defi}\label{def:T} The elements $T_m\in\essbAWeb[N](m)$ are recursively defined via:
\begin{itemize} 
\item $T_1=\id_1$,
\item $T_2=\frac{1}{N}\sum_{k=0}^{N-1} \wrap^{-k}\otimes \wrap^k$,
\item $T_{m+1}=(\id_{m-1}\otimes T_2)(T_m\otimes \id_1)$ for $m\geq 2$.
\end{itemize}
\end{defi}


 In the following we abbreviate the notation for standard basis elements of $V^m$ to $v_{a_1 a_2
\cdots a_m}:=v_{a_1}\otimes v_{a_2} \otimes \cdots \otimes v_{a_m}$.

\begin{theo} \label{thm:extweightproj} The element $\diagrep(T_m)$ is the
  endomorphism of $V^{\otimes m}$ projecting onto the sum of extremal weight spaces
  $\C\la v_{i\cdots i}\mid i\in \{1,\dots,N\}\ra$ in $\Sym^m(V)\subset V^{\otimes
  m}$.
\end{theo}
\begin{proof} For $m=1$ this is tautological. For $m=2$ we compute $\phi( \wrap^{-1}\otimes \id_1)(v_{i j})=\zeta^{-i} v_{i j}$ and $\phi(\id_1 \otimes \wrap)(v_{i j})=\zeta^{j} v_{i j}$. Thus we have:
\[\phi(T_2)(v_{i j}) =  \frac{1}{N}\sum_{k=0}^{N-1} \zeta^{k(j-i)} v_{i j} = \begin{cases} v_{i i} &\text{ if } i=j\\
0 &\text{ if } i \neq j
\end{cases}\]
For the vanishing, recall that $0=(X^N-1)=(X-1)(1+X+\cdots+X^{N-1})$, so if $X^N=1$ but $X\neq 1$, then $X$ is a zero of the cyclotomic polynomial. %This is precisely the case for $\zeta^{(j-i)}$ for $i\neq j$.
In particular, this holds for $X=\zeta^{(j-i)}$ when $i\neq j$.

For the induction step, we see immediately from the recursion that $\phi(T_{m+1})$ annihilates $v_{\epsilon_1 \epsilon_2\cdots \epsilon_{m+1}}$ unless $\epsilon_1=\cdots=\epsilon_{m}=:\epsilon$. In the remaining cases we have:
\begin{align*}
\phi(\id_{m-1}\otimes T_2)\phi(T_m\otimes \id_1)(v_{\epsilon\cdots \epsilon k}) &= \phi(\id_{m-1}\otimes T_2)(v_{\epsilon\cdots \epsilon k}) =\begin{cases} v_{\epsilon \cdots \epsilon \epsilon} &\text{ if } k=\epsilon\\ 0 &\text{ if } k\neq \epsilon \end{cases}
\end{align*} 
So $\phi(T_{m+1})$ is the extremal weight projector.
\end{proof}

In the following we show that the $T_m$ are idempotents that satisfy a number of properties analogous to the extremal weight projectors $\phi(T_m)$. This will lead to a proof that $\phi$ is indeed faithful on $\essAWeb[N]$. We start by studying properties of $T_2$.

\begin{lemma} \label{lem_T2idem}The endomorphism $T_2$ is an idempotent in $\essbAWeb[N]$ and thus also in the quotient $\essAWeb[N]$.
\end{lemma}
\begin{proof}
We compute: 
\begin{align*}T_2^2 =  \frac{1}{N^2}\sum_{k=0}^{N-1}\sum_{l=0}^{N-1} \wrap^{-k-l} \otimes \wrap^{k+l} = \frac{1}{N}\sum_{l=0}^{N-1} \wrap^{-l} \otimes \wrap^l =T_2
\end{align*}
We are using that $\wrap^{-N-i}\otimes\wrap^{N+i} = (-1)^{2N-2}{c_N}^{-1}c_N (\wrap^{-i}\otimes\wrap^{i}) =\wrap^{-i}\otimes\wrap^{i}$.
\end{proof}

\begin{lemma} \label{lem:T2comm}
In $\essbAWeb[N](3)$ we have that $\id_1 \otimes T_2$ and $T_2\otimes \id_1$ commute. 
\end{lemma}
\begin{proof}
$(\id_1 \otimes T_2)(T_2\otimes \id_1)= \frac{1}{N^2}\sum_{k,l=0}^{N-1}\wrap^{-k} \otimes \wrap^{k-l}\otimes \wrap^{l} = (\id_1 \otimes T_2)(T_2\otimes \id_1)$
\end{proof}

The following lemma states a relation that is trivially satisfied on the
representation-theoretic side, i.e. after applying $\phi$, but which is
non-obvious in $\essbAWeb[N]$ and $\essAWeb[N]$. For the latter, the result can
be deduced from \cite[Equation (37)]{CK_ann}.

\begin{lemma}\label{lem:crossabs} The idempotent $T_2$ absorbs the crossing
$s=s_1$ between its two strands. More precisely $s T_2 = T_2 s = T_2$  in
$\essbAWeb[N](2)$ and thus also in $\essAWeb[N](2)$.
\end{lemma}

The diagrammatic proof is involved and we postpone it until
Section~\ref{sec:proofs}. The reason for the result to hold after applying
$\phi$ is because $\phi(T_2)$ projects onto the span of the vectors $v_{i,i}$
(see Theorem~\ref{thm:extweightproj}), and a crossing acts as the identity on
this subspace. 



\begin{lemma} In $\essbAWeb[N](2)$ we have $\wrapi T_2 \wrap=T_2$.
\end{lemma}
\begin{proof} 
For this we rewrite $T_2$ in terms of $\wrapi\otimes \id_1= \wrapi s$ and $\id_1 \otimes \wrap= \wrap s$. Using Lemma~\ref{lem:crossabs} and $(\wrapi s)^N (\wrap s)^N = \wrap^{-N}\otimes \wrap^N= 1$, we compute:
\[\wrapi T_2 \wrap= \wrapi s T_2 \wrap= \frac{1}{N}\sum_{k=0}^{N-1} \wrapi s ( \wrapi s)^{k}(\wrap s)^k\wrap s s=T_2 s= T_2. \qedhere\]
\end{proof}

We can now state and prove the following theorem, establishing that the $T_m$'s
are indeed idempotents, and that they satisfy desirable properties with respect
to crossings, turnbacks, and wrapping morphisms.

\begin{theo} \label{thm:Tm}
  The elements $T_m$ of $\essbAWeb[N]$ satisfy the following properties: 
  \begin{enumerate}
  \item $T_m^2=T_m$; \label{step:thm:Tm:1}
  \item $T_m (\id_{k}\otimes T_n\otimes \id_{m-n-k}) =(\id_{k}\otimes T_n\otimes \id_{m-n-k}) T_m = T_m$ for $1\leq n< m$ and $0\leq k\leq m-n$; \label{step:thm:Tm:2}
  \item $(T_k\otimes \id_{m-k})(\id_{m-l}\otimes T_l)= (\id_{m-l}\otimes T_l) (T_k\otimes \id_{m-k})  = T_{m}$ for $k+l> m$; \label{item:overlap} 
\item $T_m s_i= s_i T_m=T_m$ for $m\geq 2$; \label{step:thm:Tm:4}
\item $T_m u_i= u_i T_m=0$ for $m\geq 2$; \label{step:thm:Tm:5}
 \item $\wrapi T_m \wrap=T_m$. \label{step:thm:Tm:6}
  \end{enumerate}
Here, $s_i$ and $u_i$ again refer to crossings and dumbbell webs between the strands in position $i$ and $i+1$, see Definition~\ref{def:si}.
\end{theo}
\begin{proof} We have already checked in Lemma~\ref{lem_T2idem} that $T_2$ is idempotent. From the definition, it is clear that $T_m$ is a product of $m-1$ distinct factors of the form $\id_{k}\otimes T_2 \otimes \id_{m-2-k}$ for $0\leq k\leq m-2$. Lemma~\ref{lem:T2comm} implies that these factors commute and so $T_m$ is an idempotent \eqref{step:thm:Tm:1} that absorbs smaller $T_n$, i.e. \eqref{step:thm:Tm:2}. It is also clear that overlapping projectors $T_l$ and $T_k$ combine as in \eqref{item:overlap}. The crossing absorption property of $T_2$ now implies the one for $T_m$ and crossings $s_i$ for $1\leq i \leq m-1$. 

Using crossing absorption, we obtain the rotation conjugation invariance \eqref{step:thm:Tm:6} from the $T_2$  case:
\begin{align*}
\wrapi T_m \wrap &=\wrapi (T_{m-1}\otimes \id_1)(\id_{m-2}\otimes T_2)(T_{m-1}\otimes \id_1) \wrap \\
&= (\id_1 \otimes T_{m-1})\wrapi s_{m-1}\cdots s_{2}s_{1}(T_2\otimes \id_{m-2})s_{1}s_{2}\cdots s_{m-1} \wrap (\id_1 \otimes T_{m-1})
\\
&= (\id_1 \otimes T_{m-1})((\wrapi T_2 \wrap)\otimes \id_{m-2}) (\id_1 \otimes T_{m-1}) 
\\&= (\id_1 \otimes T_{m-1})( T_2 \otimes \id_{m-2}) (\id_1 \otimes T_{m-1}) = T_{m} 
\end{align*}

This implies the missing crossing absorption relation \eqref{step:thm:Tm:4}
\[T_m s_m= T_m \wrapi s_{m-1} \wrap = \wrapi T_m  s_{m-1} \wrap = \wrapi T_m  \wrap =T_m .\] 
Finally, the $u_i$ annihilation property \eqref{step:thm:Tm:5} is equivalent to $s_i$ crossing absorption.
\end{proof}

Now we can give an alternative recursion relation for $T_m$ for $m\geq 3$. This is the direct generalization of the defining recursive relation in \cite[Definition 15]{QW}, and it is reminiscent of the Jones-Wenzl projectors.

\begin{coro} The idempotents $T_m$ satisfy the following recursion for $m\geq 3$:
\[T_m= (T_{m-1}\otimes \id_1)s_{m-1}(T_{m-1}\otimes \id_1)\]
Graphically, we write this as: 
\begin{equation}
\label{eq:recursive}
 \begin{tikzpicture}[anchorbase, scale=.3]
%% stuff inside -- inner
\draw[thick] (0,0) circle (4);
\fill[black,opacity=.2] (0,0) circle (4);
\draw[thick,fill=white] (0,0) circle (2);
\draw[dotted] (-1.05,1.05) to [out=45,in=180] (0,1.5) to [out=0,in=135] (1.05,1.05);
%% stuff inside -- outer
\draw [thick] (0,0) circle (2);
\draw[dotted] (-3.16,3.16) to [out=45,in=180] (0,4.5) to [out=0,in=135] (3.16,3.16);
%% Inner boundary circle
\draw (0,0) circle (1);
%% Outer boundary circle
\draw (0,0) circle (5);
%% web edges
\draw [very thick] (.8,.6) to (1.6,1.2);
\draw [very thick] (1,0) to (2,0);
\draw [very thick,->] (4,0) to (5,0);
\draw [very thick] (-.8,.6) to (-1.6,1.2);
\draw [very thick,->] (3.2,2.4) to (4,3);
\draw [very thick,->] (-3.2,2.4) to (-4,3);
% boundary markings
\node at (0,-1) {$*$};
\node at (0,-5) {$*$};
\draw [dashed] (0,-1) to (0,-5);
% T_{m+1}
\node at (0,2.95) {\small$T_{m}$};
\end{tikzpicture}
\;\; :=\;\;
 \begin{tikzpicture}[anchorbase, scale=.3]
%% stuff inside -- inner
\draw[thick] (0,0) circle (2.5);
\fill[black,opacity=.2] (0,0) circle (2.5);
\draw[thick,fill=white] (0,0) circle (1.5);
\draw[dotted] (-1.93,1.93) to [out=45,in=180] (0,2.75) to [out=0,in=135] (1.93,1.93);
\draw[dotted] (-0.88,0.88) to [out=45,in=180] (0,1.25) to [out=0,in=135] (0.88,0.88);
%% stuff inside -- outer
\draw[thick] (0,0) circle (4.5);
\fill[black,opacity=.2] (4.5,0) arc (0:360:4.5) -- (3.5,0) arc (360:0:3.5);
\draw [thick] (0,0) circle (3.5);
\draw[dotted] (-3.4,3.4) to [out=45,in=180] (0,4.75) to [out=0,in=135] (3.4,3.4);
\draw[dotted] (-2.29,2.29) to [out=45,in=180] (0,3.25) to [out=0,in=135] (2.29,2.29);
%% Inner boundary circle
\draw (0,0) circle (1);
%% Outer boundary circle
\draw (0,0) circle (5);
%% web edges
\draw [very thick] (.8,.6) to (1.2,.9);
\draw [very thick] (-.8,.6) to (-1.2,.9);
\draw [very thick] (-2,1.5) to (-2.8,2.1);
\fill [white] (1.4,.2) -- (2.6,.2) -- (2.6,-.2) -- (1.4,-.2);
\fill [white] (3.4,.2) -- (4.6,.2) -- (4.6,-.2) -- (3.4,-.2);
\draw [draw =white, double=black, thick, double distance=1.25pt] (1,0) -- (2.5,0) to [out=0,in=210] (2.8,2.1);
\draw [very thick,->] (2,1.5) to [out=30,in=180] (3.5,0) to (5,0);
\draw [very thick] (2.8,2.1) -- (2.8,2.1);
\draw [very thick,->] (3.6,2.7) to (4,3);
\draw [very thick,->] (-3.6,2.7) to (-4,3);
% boundary markings
\node at (0,-1) {$*$};
\node at (0,-5) {$*$};
\draw [dashed] (0,-1) to (0,-5);
% T_m
\node at (.1,1.92) {\tiny$T_{m-1}$};
\node at (0,3.95) {\tiny$T_{m-1}$};
\end{tikzpicture}
\end{equation}
\end{coro}
\begin{proof} We check this identity as follows.
\begin{align*}
(T_{m-1}\otimes \id_1)s_{m-1}(T_{m-1}\otimes \id_1) &= (T_{m-1}\otimes \id_1)s_{m-2}\cdots s_{2}s_{1}(\id_1\otimes T_{m-1}) s_1 s_2 \cdots s_{m-1}\\ 
                                                    &= (T_{m-1}\otimes \id_1)(\id_1\otimes T_{m-1} ) s_1 s_2 \cdots s_{m-1}\\
  &= T_m s_1 s_2 \cdots s_{m-1} =T_m
\end{align*}
The first equation holds by isotopy, the third by item \eqref{item:overlap} of Theorem~\ref{thm:Tm}, and the others follows from crossing absorption.
\end{proof}


\begin{lemma}\label{lem:linkedproj} Let $m,n\in \N$ with $m+n\geq 3$, then $(T_m\otimes T_n)s_m(T_m\otimes T_n)=T_{m+n}$. This means, crossing-connected projectors can be combined. 
\end{lemma}
\begin{proof} We may assume that $m\geq 2$ and compute:
  \begin{align*}
    &(\id_m\otimes T_n)(T_m\otimes \id_n)s_m(T_m\otimes \id_n)(\id_m\otimes T_n) \\
    & = (\id_m\otimes T_n)(T_{m+1}\otimes \id_{n-1})(\id_m\otimes T_n) \\
                                                                               & = T_{m+n}
  \end{align*}
  Here we have used the projector recursion \eqref{eq:recursive} and the fact that overlapping projectors can be combined, i.e. (\ref{item:overlap}) in Theorem \ref{thm:Tm}.
\end{proof}

Next we consider the images of the idempotents $T_m$ in the quotient category
$\essAWeb[N]$. Recall the morphisms $\{P_a\}_{a\in \{1,\dots, N\}}$ that were
introduced just before Theorem \ref{thm:fullness} as diagrammatic versions of
projectors on eigenspaces. We will see in Lemma \ref{lem:projexpand} that they
can be combined to give an alternate definition of the extremal weight
projectors, which amounts to saying that in the quotient category $\essAWeb[N]$,
the sum of extremal weight spaces can be broken into individual weight spaces.

\begin{lemma}\label{lem:extrabs} In $\essAWeb[N]$ we have $(P_a\otimes P_b)\circ T_2 = T_2\circ (P_a\otimes P_b) = \delta_{a,b}(P_a\otimes P_b)$.
\end{lemma}
\begin{proof} The $\C$-algebra $R:=\C[X^{\pm 1},Y^{\pm 1}]/\la X^N-1,Y^N-1\ra$ surjects onto the subalgebra of $\essAWeb[N](2)$ generated by wraps and their inverses via the map $1 \mapsto \id_2$, $X\mapsto \wrap\otimes \id_1=sD$ and $Y\mapsto \id_1\otimes \wrap=Ds$. We will check the desired equalities in $R$, where $T_2$ is represented by $\frac{1}{N}\sum_{k=0}^{N-1} X^{-k}Y^k$ and $P_a\otimes P_b$ is represented by $P_a(X)P_b(Y)$, which then implies that these equalities also hold in $\essAWeb[N](2)$.

Note that the idempotents $P_a(X)P_b(Y)$ decompose $R\cong \bigoplus_{a,b} P_a(X)P_b(Y) R$ into $1$-dimensional summands, which precisely consist of simultaneous eigenvectors for multiplication by $X$ and $Y$ with eigenvalues $\zeta^a$ and $\zeta^b$ respectively. Thus we can compute the action of $T_2$ on such an idempotent as:
\begin{align*}
  P_a(X)P_b(Y)T_2 &= P_a(X)P_b(Y) \frac{1}{N}\sum_{k=0}^{N-1} X^{-k}Y^k = P_a(X)P_b(Y) \frac{1}{N}\sum_{k=0}^{N-1} \zeta^{k(b-a)}\\
  &= \delta_{a,b}P_a(X)P_b(Y). \qedhere
\end{align*}

\end{proof}
\begin{coro}\label{cor:Ttwo} In $\essAWeb[N]$ we have 
 $T_2=\sum_{k=1}^N P_k\otimes P_k$.
\end{coro}


\begin{coro}\label{cor:crossabs} In $\essAWeb[N]$ we have $s(P_k\otimes P_k)=(P_k\otimes P_k)s=P_k\otimes P_k $.
\end{coro}
\begin{proof} We only consider composing with $s$ on the left:
\[s(P_k\otimes P_k)=s T_2(P_k\otimes P_k)=T_2(P_k\otimes P_k)=P_k\otimes P_k\]
Here we have used Lemma~\ref{lem:extrabs} twice and Lemma~\ref{lem:crossabs} in between.
\end{proof}

\begin{lemma}\label{lem:projexpand} In $\essAWeb[N]$ we have $T_m=\sum_{k=1}^N  P_k^{\otimes m}$.
\end{lemma}
\begin{proof} We have observed this for $m=2$ in Corollary~\ref{cor:Ttwo}. For $m=1$ this follows from the decomposition $1=\sum_{k=1}^N P_k$ in $\C[\wrap]/\la \wrap^N-1\ra$. Also, it is not hard to see that the $P_k$'s slide through crossings, e.g. $s \circ (P_k\otimes \id_1)=(\id_1 \otimes P_k) s$ in $\essAWeb[N](2)$. Now we proceed inductively for $m\geq 2$:
\begin{align*}
T_{m+1} &= (T_m\otimes \id_1)s_m(T_m\otimes \id_1) \\
&= \sum_{k=1}^N\sum_{l=1}^N (P_k\otimes\cdots \otimes P_k\otimes \id_1)s_m(P_l\otimes\cdots \otimes P_l\otimes \id_1) 
\\&= \sum_{k=1}^N (\id_{m-1} \otimes P_k\otimes \id_1) s_m(P_k\otimes\cdots \otimes P_k\otimes \id_1)
\\&= \sum_{k=1}^N s_m (P_k\otimes\cdots \otimes P_k\otimes P_k) =\sum_{k=1}^N  P_k^{\otimes m}
\end{align*}
Here we have used the orthogonality of the idempotents $P_k$ to proceed to the second line and the sliding property to proceed to the third line. The final crossing absorption follows from Corollary~\ref{cor:crossabs}. 
\end{proof}

\subsection{Faithfulness of the diagrammatic presentation}
We will now combine the previous results to prove the following theorem.

\begin{theo} \label{thm:faithfulness}
    The functor $\phi:\essAWeb[N]\mapsto \reph$ is faithful.
\end{theo}



We partition the proof of the theorem into three parts.

\begin{prop}\label{prop:faithful} The restriction of $\phi$ to the endomorphism algebra $\essAWeb[N](n)$ is injective.
\end{prop}
\begin{proof}
To see this result, we will exhibit a spanning set in $\essAWeb[N](n)$ that is sent under $\phi$ to a linear basis. Consider $\epsilon,\epsilon'\in \{1,\cdots,N\}^n$ so that $|\{i,\epsilon_i=k\}|=|\{i,\epsilon'_i=k\}|$ for all $k=1,\cdots, N$. Choose $\sigma_\epsilon^{\epsilon'}\in \mathfrak{S}_n$ to be the smallest in length so that $\epsilon'_{\sigma(i)}=\epsilon_i$. For example, it can be inductively defined by assigning to $1$ the smallest $r$ so that $\epsilon'_r=\epsilon_1$, etc. Recall the notation $P_k$ for the polynomial such that $P_k(D)\in \essAWeb[N](1)$ is the projector onto the $\zeta^k$ eigenspace of $D$. In $\essAWeb[N](n)$, denote $w_i=\id_{i-1}\otimes \wrap \otimes \id_{n-i}$ the complete wrap of the $i$-th strand. We define:
  \[
\phi_\epsilon^{\epsilon'}:=\sigma P_{\epsilon_n}(w_n)\cdots P_{\epsilon_1}(w_1).
\]
It is easy to see that the set $\{\phi_\epsilon^{\epsilon'}\}$ is linearly independent, because it is so under $\phi$.

On the other hand, we can deduce from Lemma~\ref{lem:upwardannulus} and its proof that this set spans $\essAWeb[N](n)$. Indeed, given the essential circle relations, we first deduce that any element $W\in \essAWeb[N](n)$ is made of a composition of elements from $\mathfrak{S}_n\rightarrow \essAWeb[N](n)$ and the wraps $w_i$. From there, using far-commutation and the formulas:
\[
w_is_{i-1}=s_{i-1}w_{i-1},\quad w_is_{i}=s_{i}w_{i+1},
\]
one can see that this gives an algebra epimorphism $ \langle w_i\rangle\rtimes \mathfrak{S}_n \twoheadrightarrow \essAWeb[N](n)$.

Now, by construction, the polynomials $\{P_k(w)\}_{k\in \{0,\dots,N-1\}}$ form a basis of $\C[w]/(w^N-1)$, and since the $w_i$'s commute, it follows that the elements $P_{\epsilon_n}(w_n)\cdots P_{\epsilon_1}(w_1)$ span $\langle w_i\rangle_{i\in \{1,\dots,n-1\}}\subset \essAWeb[N](n)$. Thus, any element in $\essAWeb[N](n)$ can be written as a linear combination of terms of the kind $\tilde{\sigma}\phi_{\epsilon}^{\epsilon}$ with $\tilde{\sigma}$ in the image of $\mathfrak{S}_n$. It remains to see that $\tilde{\sigma}$ can be assumed to be of minimal length, which is equivalent to saying that no two strands corresponding to the same value in $\epsilon$ cross.
Via isotopies, this reduces to the identities proven in Corollary~\ref{cor:crossabs}.


This proves that the set $\{\phi_{\epsilon}^{\epsilon'}\}$ spans $\essAWeb[N](n)$ and concludes the proof.
\end{proof}



\begin{prop}\label{prop:faithful2} The restriction of $\phi$ to any morphism space in $\essAWebp[N]$ is injective.
\end{prop}
\begin{proof}
Suppose a linear combination $\sum_i c_i W_i$ of webs $W_i$ in such a morphism
space is sent to zero under $\phi$. Then for each web in this linear
combination we pre-compose with a merge web $M$ and post-compose with a splitter
web $S$ (they depend only on the common source resp. target object of all webs
$W_i$) in order to obtain an endomorphism $\sum_i c_i S W_i M$ in $\essAWeb[N]$
whose source and target is a sequence of $1$s. Then we have $$\phi(\sum_i c_i S
W_i M)= \sum_i c_i \phi(S)\phi(W_i)\phi(M)=0$$ and
Proposition~\ref{prop:faithful} implies that $\sum_i c_i S W_i M=0$. Now we
consider the merge (split) web $M'$ ($S'$) obtained by reflecting $S$ ($M$) in a
horizontal line and reversing orientations. Then we compute $\sum_i c_i W_i= c
\sum_i c_i M' S W_i M S' = 0$ where $c\neq 0$ is a scalar resulting from opening
bigons; see the first relation in \eqref{eq:webrel}. 
\end{proof}

\begin{proof}[Proof of Theorem~\ref{thm:faithfulness}]
Since $\phi$ is a braided monoidal functor and since $\essAWeb[N]$ has duals,
the statement follows from Proposition~\ref{prop:faithful2}. To make this
argument more explicit, let $W$ be a linear combination of webs in some morphism
space of $\essAWeb[N]$, which is sent to zero under $\phi$. There exists a
composition of invertible bending and braiding operations similar to those used
in the proof of Theorem \ref{thm:fullness} that transforms $W$ into a linear
combination $W^\prime$ of webs in a morphism space as in
Proposition~\ref{prop:faithful2}, which is also sent to zero under $\phi$. The
proposition then implies $W^\prime=0$ and, by invertibility of the operations,
$W=0$.
\end{proof}


The final result of this section is best expressed in terms of Karoubi
envelopes, the definition of which we recall now.

\begin{defi}
The Karoubi completion of a category $\mathcal{C}$ is the category
$Kar(\mathcal{C})$ with objects given by pairs $(X,e)$, where $X$ is an object
of $\mathcal{C}$ and $e\in \Hom_\mathcal{C}(X,X)$ an idempotent. Morphisms
between $(X,e)$ and $(Y,f)$ are of the form $f\circ g \circ e$ with $g\in
\Hom_\mathcal{C}(X,Y)$. We will modify this classical definition in the following more specialized cases:
\begin{itemize}
  \item If $\mathcal{C}$ is linear but not yet additive, then we pass to the
  additive closure before taking the idempotent completion as above. The
  resulting  additive and linear category will also be denoted $Kar(\mathcal{C})$.
  \item If the morphism spaces of $\mathcal{C}$ furthermore admit a $\Z$-grading
  $\deg$, i.e. $\mathcal{C}$ is $\Z$-pre-graded, then in the definition of the
  Karoubi completion we only consider homogeneous idempotents. The resulting
  category is denoted $\Kar(\mathcal{C})^*$; it is again pre-graded. We reserve
  the notation $Kar(\mathcal{C})$ for the $\Z$-graded, additive, linear category
  whose objects are generated by formal grading shifts $\sh^k (X,e)$ of the
  objects $(X,e)$ in $\Kar(\mathcal{C})^*$ and morphism are required to be of degree
  zero. I.e. $g\colon \sh^k (X,e) \to \sh^l(Y,f)$ is required to satisfy
  $\deg(g)=l-k$.
\end{itemize}
\end{defi}

Note that if $\mathcal{C}$ is additive or monoidal, then the Karoubi completion
$Kar(\mathcal{C})$ inherits these structures.  

\begin{coro} \label{cor:diagrep} The functor $\phi$ induces an equivalence of additive, $\C$-linear pivotal categories: 
  \[
  \Kar(\essAWeb[N])^* \simeq \reph.
  \]
\end{coro}
\begin{proof} 
This follows from Theorems~\ref{thm:fullness} and \ref{thm:faithfulness} since
$\reph$ is already idempotent complete and any of its objects can be written as
the direct sum of $\phi$-images of idempotents in $\essAWeb[N]$ since $P_k$ is sent to the projection onto the span of $v_k$, see the
observation following Definition~\ref{def:Pk}.
\end{proof}

Let $\rephp$ denote the full subcategory of $\reph$ containing only those integral
$\mathfrak{h}$-representations whose weights have non-negative entries.


\begin{rema} The functor $\phi$ restricts to a fully faithful functor $\essAWebp[N] \to \rephp$ that induces an equivalence of $\C$-linear monoidal categories \[\Kar(\essAWebp[N])^* \simeq \rephp.\]
\end{rema}

\subsection{Proof of Lemma~\ref{lem:crossabs}}
\label{sec:proofs}
This section contains a proof of the fact that the $T_2$ projector absorbs crossings. It can be safely skipped on a first read-through.

In order to prove Lemma~\ref{lem:crossabs} we study the endomorphism algebra of the object $2$ in $\essbAWeb[N]$. For $k\geq 1$ we introduce the following notation: 

\begin{equation*}
E_k:=\begin{tikzpicture}[anchorbase, scale=.3]
    %% Inner boundary circle
    \draw (0,0) circle (1);
    %% Outer boundary circle
    \draw (0,0) circle (3);
    \draw [double,->] (0,1) to (0,3);
    %% Essential circle
    \draw [very thick,directed=.55] (0,0) circle (2);
    % boundary markings
    \node at (0,-1) {$*$};
    \node at (0,-3) {$*$};
        \node at (0.5,-2.4) {\tiny $k$};
    \draw [dashed] (0,-1) to (0,-3);
  \end{tikzpicture}
  \;,\;
B_k:=  \begin{tikzpicture}[anchorbase, scale=.3]
    %% Inner boundary circle
    \draw (0,0) circle (1);
    %% Outer boundary circle
    \draw (0,0) circle (3);
    \draw [double,directed=.65] (0,1) to [out=90,in=270] (-.8,1.8);
     \draw [double,->] (.8,1.8) to [out=90,in=270] (0,2.8) to (0,3);
    %% Essential circle
    \draw [very thick,directed=.55] (0,0) circle (2);
    % boundary markings
    \node at (0,-1) {$*$};
    \node at (0,-3) {$*$};
            \node at (0.5,-2.4) {\tiny $k$};
    \draw [dashed] (0,-1) to (0,-3);
  \end{tikzpicture}
    \;,\;
    A_k := \begin{tikzpicture}[anchorbase, scale=.3]
    %% Inner boundary circle
    \draw (0,0) circle (1);
    %% Outer boundary circle
    \draw (0,0) circle (3);
    \draw [double,directed=.65] (0,1) to [out=90,in=270] (.8,1.8);
     \draw [double,->] (-.8,1.8) to [out=90,in=270] (0,2.8) to (0,3);
    %% Essential circle
    \draw [very thick,directed=.55,directed=.25] (0,0) circle (2);
    % boundary markings
    \node at (0,-1) {$*$};
    \node at (0,-3) {$*$};
            \node at (.8,-2.4) {\tiny $k$-$2$};
    \draw [dashed] (0,-1) to (0,-3);
  \end{tikzpicture}
  \;,\;
  \wrap_2:=\begin{tikzpicture}[anchorbase, scale=.3]
%% Inner boundary circle
\draw (0,0) circle (1);
%% Outer boundary circle
\draw (0,0) circle (3);
%% wrapping around
\draw [double,->] (0,1) .. controls (0,1.7) and (-1.6,1) .. (-1.6,0) to [out=-90,in=180] (0,-1.8) to [out=0,in=-90] (2.1,0) to [out=90,in=-90] (0,3); 
% boundary markings
\node at (0,-1) {$*$};
\node at (0,-3) {$*$};
\draw [dashed] (0,-1) to (0,-3);
\end{tikzpicture}
\end{equation*}
Here we set $A_1=0$, $B_2=\wrap_2$ and $A_2=\id$, and doubled edges stand for 2-labeled edges. Note that in the definition of $B_k$, we haven't depicted the orientation of one of the strands: this is because it depends on $k$. More precisely, we have:
  \[
  B_1:=  \begin{tikzpicture}[anchorbase, scale=.3]
    %% Inner boundary circle
    \draw (0,0) circle (1);
    %% Outer boundary circle
    \draw (0,0) circle (3);
    \draw [double,directed=.65] (0,1) to [out=90,in=270] (-.8,1.8);
     \draw [double,->] (.8,1.8) to [out=90,in=270] (0,2.8) to (0,3);
    %% Essential circle
    \draw [very thick,directed=.55,rdirected=.26] (0,0) circle (2);
    % boundary markings
    \node at (0,-1) {$*$};
    \node at (0,-3) {$*$};
    \node at (0.5,-2.4) {\tiny $1$};
    \node at (0.3,1.5) {\tiny $1$};
    \draw [dashed] (0,-1) to (0,-3);
  \end{tikzpicture}\quad,\quad
  B_2=\wrap_2,\quad
  B_k:=  \begin{tikzpicture}[anchorbase, scale=.3]
    %% Inner boundary circle
    \draw (0,0) circle (1);
    %% Outer boundary circle
    \draw (0,0) circle (3);
    \draw [double,directed=.65] (0,1) to [out=90,in=270] (-.8,1.8);
     \draw [double,->] (.8,1.8) to [out=90,in=270] (0,2.8) to (0,3);
    %% Essential circle
    \draw [very thick,directed=.55,directed=.26] (0,0) circle (2);
    % boundary markings
    \node at (0,-1) {$*$};
    \node at (0,-3) {$*$};
    \node at (0.5,-2.4) {\tiny $k$};
    \draw [dashed] (0,-1) to (0,-3);
  \end{tikzpicture}
  \quad\text{if}\; k\geq 3.
\]
 Note also that $A_k=B_k=E_k=0$ for $k>N$.

\begin{lemma}\label{lem:endtwo} The following statements hold in the endomorphism algebra of the $2$-labeled upward point in $\essbAWeb[N]$:
\begin{enumerate}
\item $E_k = \delta_{k,N} c_N$, \label{step:lem:endtwo:1}
\item $\wrap_2$ is invertible and central, \label{step:lem:endtwo:2}
\item $B_k= A_k \wrap_2$ for $k\geq 2$, \label{step:lem:endtwo:3}
\item $B_1 A_k = - E_{k-1} + A_{k-1} \wrap_2 + A_{k+1}$ for $k\geq 2$, \label{step:lem:endtwo:4}
\item $B_N=E_N=c_N$ and thus $A_{N}=c_N \wrap_2^{-1}$. \label{step:lem:endtwo:5}
\end{enumerate}
\end{lemma}
Note that only \eqref{step:lem:endtwo:1} and \eqref{step:lem:endtwo:5} depend on the value of $N$.
\begin{proof}
  \eqref{step:lem:endtwo:1} holds by definition of $\essbAWeb[N]$, \eqref{step:lem:endtwo:2} and \eqref{step:lem:endtwo:3} follow from isotopies. For \eqref{step:lem:endtwo:4} we resolve the crossing in $E_{k-1}$ to obtain:
  \begin{equation*}
\begin{tikzpicture}[anchorbase, scale=.25]
    %% Inner boundary circle
    \draw (0,0) circle (1);
    %% Outer boundary circle
    \draw (0,0) circle (5);
    \draw [double,->] (0,1) to (0,5);
    %% Essential circle
    \draw [very thick,directed=.55] (0,0) circle (3);
    % boundary markings
    \node at (0,-1) {$*$};
    \node at (0,-5) {$*$};
    \draw [dashed] (0,-1) to (0,-5);
    \node at (1.2,-3.6) {\tiny $k-1$};
  \end{tikzpicture}
  \;\; = \;\;
  \begin{tikzpicture}[anchorbase, scale=.25]
    %% Inner boundary circle
    \draw (0,0) circle (1);
    %% Outer boundary circle
    \draw (0,0) circle (5);
     \draw [very thick,directed=.55] (0,0) circle (3);
    \draw [double,directed=.65] (0,1) to [out=90,in=270] (-1.2,2.7);
     \draw [double,->] (1.2,2.7) to [out=90,in=270] (0,4.5) to (0,5);
    \node at (0,-1) {$*$};
    \node at (0,-5) {$*$};
    \draw [dashed] (0,-1) to (0,-5);
    \node at (1.2,-3.6) {\tiny $k-1$};
  \end{tikzpicture}
 \;\; -\;\;
   \begin{tikzpicture}[anchorbase, scale=.25]
    %% Inner boundary circle
    \draw (0,0) circle (1);
    %% Outer boundary circle
    \draw (0,0) circle (5);
    % \draw [very thick,directed=.55] (0,0) circle (3);
     \draw[very thick,directed=.11] (180:3) arc (180:0+360:3);
      \draw [very thick] (3,0) to [out=90,in=0](1,3);
      \draw [very thick] (-1,3) to [out=180,in=90](-3,0);
    \draw [double,directed=.65] (0,1) to (0,2);
    \draw[very thick, directed=.55] (0,2) to (-1,3);
    \draw[very thick, directed=.55] (1,3) to (0,2);
    \draw[very thick, directed=.55] (-1,3) to (0,4);
    \draw[very thick, directed=.55] (1,3) to (0,4);
     \draw [double,->] (0,4) to (0,5);
    \node at (0,-1) {$*$};
    \node at (0,-5) {$*$};
    \draw [dashed] (0,-1) to (0,-5);
    \node at (1.2,-3.6) {\tiny $k-1$};
  \end{tikzpicture}        
\;\;+ \;\;
  \begin{tikzpicture}[anchorbase, scale=.25]
    %% Inner boundary circle
    \draw (0,0) circle (1);
    %% Outer boundary circle
    \draw (0,0) circle (5);
     \draw [very thick,directed=.55] (0,0) circle (3);
    \draw [double,directed=.65] (0,1) to [out=90,in=270] (1.2,2.7);
     \draw [double,->] (-1.2,2.7) to [out=90,in=270] (0,4.5) to (0,5);
    \node at (0,-1) {$*$};
    \node at (0,-5) {$*$};
    \draw [dashed] (0,-1) to (0,-5);
    \node at (1.2,-3.6) {\tiny $k-1$};
  \end{tikzpicture}        
\end{equation*}
Here the first and third summands are $B_{k-1}=A_{k-1}\wrap_2$ and $A_{k+1}$ respectively. The web in the second summand simplifies as follows: 
\[   \begin{tikzpicture}[anchorbase, scale=.25]
    %% Inner boundary circle
    \draw (0,0) circle (1);
    %% Outer boundary circle
    \draw (0,0) circle (5);
     \draw[very thick,directed=.11] (180:3) arc (180:0+360:3);
      \draw [very thick] (3,0) to [out=90,in=0](1,3);
      \draw [very thick] (-1,3) to [out=180,in=90](-3,0);
    \draw [double,directed=.65] (0,1) to (0,2);
    \draw[very thick, directed=.55] (0,2) to (-1,3);
    \draw[very thick, directed=.55] (1,3) to (0,2);
    \draw[very thick, directed=.55] (-1,3) to (0,4);
    \draw[very thick, directed=.55] (1,3) to (0,4);
     \draw [double,->] (0,4) to (0,5);
    \node at (0,-1) {$*$};
    \node at (0,-5) {$*$};
    \draw [dashed] (0,-1) to (0,-5);
    \node at (1.2,-3.6) {\tiny $k-1$};
  \end{tikzpicture}        
  \;\;=\;\;
   \begin{tikzpicture}[anchorbase, scale=.25]
    %% Inner boundary circle
    \draw (0,0) circle (1);
    %% Outer boundary circle
    \draw (0,0) circle (5);
     \draw[very thick,directed=.11] (180:2) arc (180:0+360:2);
     \draw[very thick,directed=.11] (180:4) arc (180:0+360:4);
      \draw [very thick] (2,0) to [out=90,in=0](0,2);
      \draw [very thick] (4,0) to [out=90,in=0](0,4);
      \draw [very thick, directed=.55] (-1,3) to [out=180,in=90](-3,1);
      \draw [very thick] (-4,0) to [out=90,in=225](-3,1);
      \draw [very thick] (-2,0) to [out=90,in=315](-3,1);
    \draw [double,directed=.65] (0,1) to (0,2);
    \draw[very thick, directed=.55] (0,2) to (-1,3);
    \draw[very thick, directed=.55] (-1,3) to (0,4);
     \draw [double,->] (0,4) to (0,5);
    \node at (0,-1) {$*$};
    \node at (0,-5) {$*$};
    \draw [dashed] (0,-1) to (0,-5);
    \node at (1.2,-2.6) {\tiny $k-2$};
    \node at (1.2,-4.5) {\tiny $1$};
  \end{tikzpicture}    
  \;\;=\;\;
    \begin{tikzpicture}[anchorbase, scale=.25]
    %% Inner boundary circle
    \draw (0,0) circle (1);
    %% Outer boundary circle
    \draw (0,0) circle (5);
     \draw [very thick,directed=.55, rdirected=.32] (0,0) circle (4);
     \draw [very thick,directed=.55, directed=.34] (0,0) circle (2);
    \draw [double,directed=.65] (0,1) to (0,2);
     \draw [double,->] (-1.42,1.42) to  (-2.84,2.84);
     \draw [double,->] (0,4) to (0,5);
    \node at (0,-1) {$*$};
    \node at (0,-5) {$*$};
    \draw [dashed] (0,-1) to (0,-5);
  \node at (1.2,-2.6) {\tiny $k-2$};
    \node at (1.2,-4.5) {\tiny $1$};
  \end{tikzpicture}
  \;\; =\;\; B_1 A_{k}.
  \qedhere \]
\end{proof}

As a corollary, we get that the elements $A_k$ can be written in terms of powers of $B_1$ and $D_2$:
\begin{coro} \label{cor:endtwob}	  The elements $A_k$ for $k<N+2$ satisfy the recursion $A_k := B_1 A_{k-1} -A_{k-2}\wrap_2$ for $k\geq 5$ with initial conditions $A_3=B_1$ and $A_4=B_1^2-\wrap_2$.
\end{coro}
\begin{proof} We induct on $k$. For $k=1<N$ we use Lemma~\ref{lem:endtwo} to obtain $B_1=B_1A_2 = -E_1 +A_1 \wrap_2 +A_3=A_3$. Similarly, for $k=2<N$, we get from Lemma~\ref{lem:endtwo} that $B_1^2 - \wrap_2= B_1 A_3-A_2\wrap_2 = -E_2 + A_4=A_4$. For the recurrence relation we compute:
  $B_1 A_{k-1} -A_{k-2}\wrap_2 = -E_{k-2} + A_{k} = A_{k}$ if $5\leq k<N+2$.
\end{proof}

\begin{rema} We will now find expressions for $T_2$ which are more convenient in the following proof.
First we introduce the notation $t=\wrapi s \wrap$ for the rotation conjugate of the crossing. Now we use the facts that $\wrap^2$ is central in $\essbAWeb[N](2)$ and that $(s\wrap)^{-1}$ and $(\wrap s)$ commute to write $(s \wrap)^{-k} (\wrap s)^k = (t s)^k$, which gives: 
\[T_2=\frac{1}{N}\sum_{k=0}^{N-1} (t s)^k\]  
Further, we resolve the crossings as $s=\id_2-u$ and $t=\id_2-v$ where $u=u_1$ is the 2-labeled dumbbell and $v$ its conjugate. So we have:
\begin{equation}
\label{eqn_Ttwob}
T_2=\frac{1}{N}\sum_{k=0}^{N-1} ((\id_2-v)(\id_2-u))^k
\end{equation}
\end{rema}

If we write $M$ and $S$ for the merge and split vertices on two strands, such that $SM=u$, we get the following equality for $k\geq 2$:
\[\underbrace{\cdots v u v u}_{k \text{ factors}} = \wrap^{1-k} S B_1^{k-1} M\]  

We will now attempt to rewrite the expressions $X_n:=\wrap^{1-n} S A_{n+1} M$ in terms of powers of $B_1$. We can then use the relation $X_N=0$ to deduce a relation between compositions of the webs $u$ and $v$. 
To this end we define $R_{2k-1}:=\underbrace{(\id_2-v)(\id_2-u)\cdots (\id_2-v)}_{2k-1 \text{ factors}}$ and $S_x:=\sum_{k=1}^x R_{2k-1}$. 

\begin{lemma} \label{lem:XY}For $2\leq n \leq N$ we have:
\[X_n = \begin{cases} 
u - R_{n-1}u - u S_{n/2-1} u & \text{ even } n \\
u - u S_{(n-1)/2} u & \text{ odd } n
\end{cases}\]
\end{lemma}
\begin{proof} We will use the notation $Y_n$ for the entries on the right-hand side of the equation in the statement of the lemma. The proof of $X_n=Y_n$ proceeds by induction on $n$. For $n=2$ we have $X_2=v u = u-(\id_2-v)u = Y_2$. Similarly, for $n=3$, we have
\[ X_3=\wrap^{-2} S A_{4} M = \wrap^{-2} S (B_1^2-\wrap_2) M = uvu-u = u - u(1-v)u = Y_3\] since $u^2=2 u$ by the bigon relation. We prove the remaining cases recursively. For this, note that the elements $X_n=\wrap^{1-n} S A_{n+1} M$ inherit a recurrence relation from the elements $A_{n+1}$:
\begin{align}\label{eqn:Xrec}
X_n=\wrap^{1-n} S A_{n+1} M &= \wrap^{1-n} S B_1 A_{n} M - \wrap^{1-n} S A_{n-1} \wrap_2 M \\
\nonumber
&= \wrap^{1-n} S B_1 A_{n} M - \wrap^{3-n} S A_{n-1} M  \\
\nonumber
&= \begin{cases}
v X_{n-1} - X_{n-2} & \text{even } n\\
u X_{n-1} - X_{n-2} & \text{odd } n
\end{cases} 
\end{align}
Here we have used $\wrap^{1-n}SB_1=v \wrap^{2-n}S$ for even $n$ and $\wrap^{1-n}SB_1=u \wrap^{2-n} S$ for odd $n$. Now it remains to check that the $Y_n$ also satisfy this recurrence \eqref{eqn:Xrec}. Indeed, for odd $N>4$ we can verify:
\begin{align*}
  Y_n-uY_{n-1}+Y_{n-2} &= (u - u S_{(n-1)/2} u) - u(u - R_{n-2}u - u S_{(n-3)/2} u) \\
  &\quad + (u - u S_{(n-3)/2} u)\\
&=  u (-S_{(n-1)/2} + R_{n-2} + S_{(n-3)/2}) u =0
\end{align*}
Here we have used that $S_x-S_{x-1}= R_{2x-1}$. In order to check the recurrence for $Y_n$ in the case of even $n$ we need an auxiliary computation. For odd $x\geq 1$ we have
\begin{align*}S_x  &= (\id_2-v)  + (\id_2-v)(\id_2-u) S_{x-1}\\
  &= (\id_2-v)  + (\id_2-v)S_{x-1}  - u S_{x-1} + v u S_{x-1}  
\\
 &= 2\id_2 -v  + (\id_2-u)S_{x-2}  - u S_{x-1} + v u S_{x-1}  
\end{align*} which implies:
\begin{align*} v u S_{n/2-1} u &= - 2u + v u +S_{n/2} u- (\id_2-u) S_{n/2-2} u +u S_{n/2-1}u\\
&= 
-2 u + v u + R_{n-1} u+ R_{n-3}u + u S_{n/2-2} u +u S_{n/2-1}u
\end{align*}
Here we have used $(\id_2-v)S_x = \id_2 + (\id_2-u)S_{x-1}$. Now we check the recurrence for even $n>3$:
\begin{align*}
  Y_n-v Y_{n-1}+Y_{n-2} &= (u - R_{n-1}u - u S_{n/2-1} u) - v (u - u S_{n/2-1} u) \\
  &\quad + (u - R_{n-3}u - u S_{n/2-2} u) \\
  &= 2u -v u - R_{n-1}u -R_{n-3}u - u S_{n/2-2} u- u S_{n/2-1} u  \\
  &\quad + v u S_{n/2-1} u \\
  &= 0 
\end{align*}
This completes the proof of the Lemma.
\end{proof}

\begin{proof}[Proof of Lemma~\ref{lem:crossabs}]
We only prove $T_2 s=T_2$, which is equivalent to $N T_2 u=0$ by expanding the crossing and multiplying by $N$. Using \eqref{eqn_Ttwob} we compute: 
\begin{align*} 
  N T_2 u &=  \sum_{k=0}^{N-1} ((\id_2-v)(\id_2-u))^k u = u - \sum_{k=1}^{N-1} ((\id_2-v)(\id_2-u))^{k-1} (\id_2-v)u \\
  &= u - \sum_{k=1}^{N-1} R_{2k -1} u
\end{align*}
Now note that $\id_2 = (s \wrap)^{-N} (\wrap s)^N = ((\id_2-v)(\id_2-u))^N$ implies that 
\begin{equation}
\label{eqn:rflip}
R_{2k-1} = \left( (\id_2-u) R_{2N-2k-1} (\id_2-u)\right)^{-1} = (\id_2-u) R_{2N-2k-1} (\id_2-u).
\end{equation}
Now we distinguish two cases. For even $N$ we expand:
\begin{align*} 
N T_2 u  &= u - \sum_{k=1}^{N/2} R_{2k -1} u - \sum_{k=N/2+1}^{N-1} R_{2k -1} u
\\
&= u - \sum_{k=1}^{N/2} R_{2k -1} u - (\id_2-u)\sum_{l=1}^{N/2-1} R_{2l -1}(\id_2-u)u\\
&= u - \sum_{k=1}^{N/2} R_{2k -1} u + (\id_2-u)\sum_{l=1}^{N/2-1} R_{2l -1}u = u - R_{N-1} u - u S_{N/2-1} u = X_N 
\end{align*}
Here we have used \eqref{eqn:rflip} for the second equality and Lemma~\ref{lem:XY} for the last equality. For odd $N$ we expand analogously:
\begin{align*} 
N T_2 u  &= u - \sum_{k=1}^{(N-1)/2} R_{2k -1} u - \!\!\!\!\sum_{k=(N+1)/2+1}^{N-1} R_{2k -1} u\\
&= u - \sum_{k=1}^{(N-1)/2} R_{2k -1} u + (\id_2-u)\sum_{l=1}^{(N-1)/2} R_{2k -1} u = u- u S_{(N-1)/2}u = X_N
\end{align*}
We conclude the proof by noting that $X_N=\wrap^{1-N} S A_{N+1} M=0$ in $\essbAWeb[N]$ since $A_{N+1}=0$. 
\end{proof}

\begin{rema}
The expression $X_N$ in the rewritten form in Lemma~\ref{lem:XY} expresses the
longest Kazhdan--Lusztig basis element $\underline{H}_{s t
s\dots}=\underline{H}_{t s t\dots}$ in the type $I_2(N)$ Hecke algebra in terms
of products of $\underline{H}_s:=u$ and $\underline{H}_t:=v$, see \cite[Section
2.3]{EW}. In particular, the relation $X^N=0$ suggests that $\essbAWeb[N](2)$ is
related to a quotient of the Hecke algebra by the $2$-cell containing the basis
element associated to the longest word. 
\end{rema}

\begin{thebibliography}{MMMMMM}
\bibitem[1]{APS} Asaeda, Marta M.; Przytycki, Józef H.; Sikora, Adam S. Categorification of the Kauffman bracket skein module of I-bundles over surfaces, Algebr. Geom. Topol., Volume 4 (2004), pp. 1177-1210 \href{https://alco.centre-mersenne.org/articles/10.5802/alco.330/#APS}{Publisher reference}.
\bibitem[2]{BN2} Bar-Natan, Dror Khovanov’s homology for tangles and cobordisms, Geom. Topol., Volume 9 (2005), pp. 1443-1499 \href{https://alco.centre-mersenne.org/articles/10.5802/alco.330/#BN2}{Publisher reference}.
\bibitem[3]{BPW} Beliakova, Anna; Putyra, Krzysztof K.; Wehrli, Stephan M. Quantum link homology via trace functor I, Invent. Math., Volume 215 (2019) no. 2, pp. 383-492 \href{https://alco.centre-mersenne.org/articles/10.5802/alco.330/#BPW}{Publisher reference}.
\bibitem[4]{Blan} Blanchet, Christian An oriented model for Khovanov homology, J. Knot Theory Ramifications, Volume 19 (2010) no. 2, pp. 291-312 \href{https://alco.centre-mersenne.org/articles/10.5802/alco.330/#Blan}{Publisher reference}.
\bibitem[5]{BW1} Bonahon, Francis; Wong, Helen Representations of the Kauffman bracket skein algebra III: closed surfaces and naturality, Quantum Topol., Volume 10 (2019) no. 2, pp. 325-398 \href{https://alco.centre-mersenne.org/articles/10.5802/alco.330/#BW1}{Publisher reference}.
\bibitem[6]{Bou} Bousseau, Pierrick Strong positivity for the skein algebras of the 4-punctured sphere and of the 1-punctured torus, Comm. Math. Phys., Volume 398 (2023) no. 1, pp. 1-58 \href{https://alco.centre-mersenne.org/articles/10.5802/alco.330/#Bou}{Publisher reference}.
\bibitem[7]{Bru} Brundan, J. Representations of the oriented skein category, 2017 \href{https://alco.centre-mersenne.org/articles/10.5802/alco.330/#Bru}{Publisher reference}.
\bibitem[8]{CK_ann} Cautis, Sabin; Kamnitzer, Joel Quantum K-theoretic geometric Satake: the  SL  n  case, Compos. Math., Volume 154 (2018) no. 2, pp. 275-327 \href{https://alco.centre-mersenne.org/articles/10.5802/alco.330/#CK_ann}{Publisher reference}.
\bibitem[9]{CKM} Cautis, Sabin; Kamnitzer, Joel; Morrison, Scott Webs and quantum skew Howe duality, Math. Ann., Volume 360 (2014) no. 1-2, pp. 351-390 \href{https://alco.centre-mersenne.org/articles/10.5802/alco.330/#CKM}{Publisher reference}.
\bibitem[10]{ETW} Ehrig, Michael; Tubbenhauer, Daniel; Wedrich, Paul Functoriality of colored link homologies, Proc. Lond. Math. Soc. (3), Volume 117 (2018) no. 5, pp. 996-1040 \href{https://alco.centre-mersenne.org/articles/10.5802/alco.330/#ETW}{Publisher reference}.
\bibitem[11]{EW} Elias, Ben; Williamson, Geordie Soergel calculus, Represent. Theory, Volume 20 (2016), pp. 295-374 \href{https://alco.centre-mersenne.org/articles/10.5802/alco.330/#EW}{Publisher reference}.
\bibitem[12]{FG} Frohman, Charles; Gelca, Răzvan Skein modules and the noncommutative torus, Trans. Amer. Math. Soc., Volume 352 (2000) no. 10, pp. 4877-4888 \href{https://alco.centre-mersenne.org/articles/10.5802/alco.330/#FG}{Publisher reference}.
\bibitem[13]{GLW} Grigsby, J. Elisenda; Licata, Anthony M.; Wehrli, Stephan M. Annular Khovanov homology and knotted Schur–Weyl representations, Compos. Math., Volume 154 (2018) no. 3, pp. 459-502 \href{https://alco.centre-mersenne.org/articles/10.5802/alco.330/#GLW}{Publisher reference}.
\bibitem[14]{Jon2} Jones, V. F. R. Index for subfactors, Invent. Math., Volume 72 (1983) no. 1, pp. 1-25 \href{https://alco.centre-mersenne.org/articles/10.5802/alco.330/#Jon2}{Publisher reference}.
\bibitem[15]{Kh1} Khovanov, Mikhail A categorification of the Jones polynomial, Duke Math. J., Volume 101 (2000) no. 3, pp. 359-426 \href{https://alco.centre-mersenne.org/articles/10.5802/alco.330/#Kh1}{Publisher reference}.
\bibitem[16]{Kh7} Khovanov, Mikhail Categorifications of the colored Jones polynomial, J. Knot Theory Ramifications, Volume 14 (2005) no. 1, pp. 111-130 \href{https://alco.centre-mersenne.org/articles/10.5802/alco.330/#Kh7}{Publisher reference}.
\bibitem[17]{KhR} Khovanov, Mikhail; Rozansky, Lev Matrix factorizations and link homology, Fund. Math., Volume 199 (2008) no. 1, pp. 1-91 \href{https://alco.centre-mersenne.org/articles/10.5802/alco.330/#KhR}{Publisher reference}.
\bibitem[18]{KSa} Khovanov, Mikhail; Sazdanovic, Radmila A categorification of one-variable polynomials, Proceedings of FPSAC 2015 (Discrete Math. Theor. Comput. Sci. Proc.), Assoc. Discrete Math. Theor. Comput. Sci., Nancy (2015), pp. 937-948 \href{https://alco.centre-mersenne.org/articles/10.5802/alco.330/#KSa}{Publisher reference}.
\bibitem[19]{Kup} Kuperberg, Greg Spiders for rank 2 Lie algebras, Comm. Math. Phys., Volume 180 (1996) no. 1, pp. 109-151 \href{https://alco.centre-mersenne.org/articles/10.5802/alco.330/#Kup}{Publisher reference}.
\bibitem[20]{LTV} Lacabanne, Abel; Tubbenhauer, Daniel; Vaz, Pedro Annular webs and Levi subalgebras, J. Comb. Algebra, Volume 7 (2023) no. 3-4, pp. 283-326 \href{https://alco.centre-mersenne.org/articles/10.5802/alco.330/#LTV}{Publisher reference}.
\bibitem[21]{Le} Lê, Thang T. Q. On positivity of Kauffman bracket skein algebras of surfaces, Int. Math. Res. Not. IMRN (2018) no. 5, pp. 1314-1328 \href{https://alco.centre-mersenne.org/articles/10.5802/alco.330/#Le}{Publisher reference}.
\bibitem[22]{MS} Morton, Hugh; Samuelson, Peter The HOMFLYPT skein algebra of the torus and the elliptic Hall algebra, Duke Math. J., Volume 166 (2017) no. 5, pp. 801-854 \href{https://alco.centre-mersenne.org/articles/10.5802/alco.330/#MS}{Publisher reference}.
\bibitem[23]{MOY} Murakami, Hitoshi; Ohtsuki, Tomotada; Yamada, Shuji Homfly polynomial via an invariant of colored plane graphs, Enseign. Math. (2), Volume 44 (1998) no. 3-4, pp. 325-360 \href{https://alco.centre-mersenne.org/articles/10.5802/alco.330/#MOY}{Publisher reference}.
\bibitem[24]{Prz1} Przytycki, Józef H. Skein modules of 3-manifolds, Bull. Polish Acad. Sci. Math., Volume 39 (1991) no. 1-2, pp. 91-100 \href{https://alco.centre-mersenne.org/articles/10.5802/alco.330/#Prz1}{Publisher reference}.
\bibitem[25]{Queff_aff} Queffelec, Hoel Skein modules from skew Howe duality and affine extensions, SIGMA Symmetry Integrability Geom. Methods Appl., Volume 11 (2015), Paper no. 030, 36 pages \href{https://alco.centre-mersenne.org/articles/10.5802/alco.330/#Queff_aff}{Publisher reference}.
\bibitem[26]{QR2} Queffelec, Hoel; Rose, David E. V. Sutured annular Khovanov-Rozansky homology, Trans. Amer. Math. Soc., Volume 370 (2018) no. 2, pp. 1285-1319 \href{https://alco.centre-mersenne.org/articles/10.5802/alco.330/#QR2}{Publisher reference}.
\bibitem[27]{QS2} Queffelec, Hoel; Sartori, Antonio Mixed quantum skew Howe duality and link invariants of type A, J. Pure Appl. Algebra, Volume 223 (2019) no. 7, pp. 2733-2779 \href{https://alco.centre-mersenne.org/articles/10.5802/alco.330/#QS2}{Publisher reference}.
\bibitem[28]{QW} Queffelec, Hoel; Wedrich, Paul Extremal weight projectors, Math. Res. Lett., Volume 25 (2018) no. 6, pp. 1911-1936 \href{https://alco.centre-mersenne.org/articles/10.5802/alco.330/#QW}{Publisher reference}.
\bibitem[29]{QW3} Queffelec, Hoel; Wedrich, Paul Khovanov homology and categorification of skein modules, Quantum Topol., Volume 12 (2021) no. 1, pp. 129-209 \href{https://alco.centre-mersenne.org/articles/10.5802/alco.330/#QW3}{Publisher reference}.
\bibitem[30]{Thu} Thurston, Dylan Paul Positive basis for surface skein algebras, Proc. Natl. Acad. Sci. USA, Volume 111 (2014) no. 27, pp. 9725-9732 \href{https://alco.centre-mersenne.org/articles/10.5802/alco.330/#Thu}{Publisher reference}.
\bibitem[31]{TVW} Tubbenhauer, Daniel; Vaz, Pedro; Wedrich, Paul Super q-Howe duality and web categories, Algebr. Geom. Topol., Volume 17 (2017) no. 6, pp. 3703-3749 \href{https://alco.centre-mersenne.org/articles/10.5802/alco.330/#TVW}{Publisher reference}.
\bibitem[32]{Tur} Turaev, Vladimir G. Skein quantization of Poisson algebras of loops on surfaces, Ann. Sci. École Norm. Sup. (4), Volume 24 (1991) no. 6, pp. 635-704 \href{https://alco.centre-mersenne.org/articles/10.5802/alco.330/#Tur}{Publisher reference}.
\bibitem[33]{Wen} Wenzl, Hans On sequences of projections, C. R. Math. Rep. Acad. Sci. Canada, Volume 9 (1987) no. 1, pp. 5-9 \href{https://alco.centre-mersenne.org/articles/10.5802/alco.330/#Wen}{Publisher reference}.
\end{thebibliography}
\end{document}
